\chapter{The Hitchin-Kobayashi correspondence} \label{chapter_HK_correspondence}
In this chapter, all the different topics we have treated come together in the Hitchin-Kobayashi correspondence. We will not treat the vector bundle version, even though it was the first version to appear, and we will jump straight to the principal bundle version. First, we give an overview of the result when the structure group is connected. We then consider the case where the structure group is non-connected, and we apply the techniques from appendix \ref{appendix_lie} to the study of this new situation. We also see how the moduli space of solutions of the hermitian Yang-Mills equation plays a role in an infinite-dimensional analogue of the Kempf-Ness theorem. To be fair, we will only cover here the correspondence for principal bundles over Riemann surfaces, which is Ramanathan's generalization of the Narasimhan-Seshadri theorem, but we will keep calling it the Hitchin-Kobayashi correspondence (HK correspondence from now on).

\section{The HK correspondence for connected groups}
For the remainder of the section, we fix a complex connected reductive Lie group \(G\) and let \(K\) be a maximal compact subgroup. Let \(\mathcal{P}_G\) be a holomorphic principal \(G\)-bundle over some compact Riemann surface \(X\). The statement of the Hitchin-Kobayashi correspondence is the following: a holomorphic principal \(G\)-bundle over a compact Kähler manifold is polystable if and only if there is a reduction of the structure group to \(K\) whose Chern connection satisfies the hermitian Yang-Mills equation. We discussed the equation in chapter \ref{chapter_HYM_equation}, and we discuss the stability condition below. We will use extensively the notions of parabolic subgroups introduced in appendix \ref{appendix_lie}. It is important to remark that one may find several definitions of stability in the literature. They are all equivalent, even though it is not obvious at first glance. Our definitions will follow \parencite{i1999hitchin} and \parencite{garcia2009hitchin}. From now on, we fix \(\langle \cdot, \cdot \rangle\colon \mathfrak{g} \times \mathfrak{g} \to \C\) a non-degenerate \(\C\)-bilinear \(\Ad{}\)-invariant pairing on \(\mathfrak{g}\) that restricts to a non-degenerate pairing on \(\mathfrak{k}\). We fix the notation \(\mathscr{G}:=\Gamma(X; P_K \times_K K)\) and \(\mathscr{G}^\C=\Gamma(X; P_G \times_G G)\). We will write \(\exp_{\mathscr{G}}\) and \(\exp_{\mathscr{G}^\C}\) for their respective exponential maps. 

\Definition{}{Let \(\mathfrak{g}\) be a complex Lie algebra. \\
\noindent
\begin{minipage}{0.77\textwidth}
    A \textit{character}\index{Character} is a complex linear map \(\chi\colon \mathfrak{g}\to \C\) that factorizes through the quotient \(q\colon \mathfrak{g}\to \mathfrak{g}/[\mathfrak{g}, \mathfrak{g}]\). That is, there is a complex linear map \(\bar{\chi}\colon \mathfrak{g}/[\mathfrak{g}, \mathfrak{g}] \to \C\) such that the diagram commutes.
\end{minipage}\hfill%
\begin{minipage}{0.22\textwidth}
        \begin{center}
    \begin{tikzcd}[row sep=normal, column sep=normal,  ampersand replacement=\&]
            \mathfrak{g} \arrow[d, "{\displaystyle q}", '] \arrow[r, "{\displaystyle \chi}"] \& \C \\
            \mathfrak{g}/[\mathfrak{g}, \mathfrak{g}] \arrow[ur, "{\displaystyle \bar{\chi}}",'] \& 
        \end{tikzcd}
\end{center}
\end{minipage}
}

The following lemma gives examples of characters on parabolic algebras of the form \(\mathfrak{p}_\zeta\) for \(\zeta \in \ii \mathfrak{k}\).

\Lemma{}{Let \(\zeta \in \mathfrak{k}\subset \mathfrak{g}\). We define \(\chi_\zeta\in \mathfrak{g}^*\) by \label{lemma_character}
\[
\chi_\zeta(\xi)=\langle \zeta, \xi \rangle \hspace{1 cm} \xi \in \mathfrak{g}. 
\]
Then, the restriction of \(\chi_\zeta\) to \(\mathfrak{p}_\zeta\) (which we will denote by the same symbol) is a character on \(\mathfrak{p}_\zeta^*\). \label{lemma_antidominant_character}
}
\begin{proof}
    \(\chi_\zeta\) is clearly \(\C\)-linear, and so is its restriction. Observe that \(\zeta \in \mathfrak{p}_\zeta\). This follows from the definition of \(\mathfrak{p}_\zeta\), and the fact that \(\Ad{\exp(t\zeta)}(\zeta)=\zeta\) for all \(t \in \R\). To check that it factors through \([\mathfrak{p}_\zeta, \mathfrak{p}_\zeta]\), we must see that \(\langle \zeta, [x, y]\rangle =0\) for \(x,y \in \mathfrak{p}_\zeta\). Observe that \(\langle \zeta, [x,y]\rangle = \langle [x,\zeta], y \rangle\). This follows from \(\Ad{}\)-invariance, and from the fact that \(\mathrm{ad}(x)(y)=[x,y]=\diff_e \Ad{}(x)(y)\) \parencite[Th. 20.27]{lee_introSM}:
    \begin{align*}
        \langle \zeta, [x,y] \rangle & = \langle \zeta, \mathrm{ad}(x)(y) \rangle = \langle \zeta, \diff_e \Ad{}(x)(y) \rangle = \bigg\langle \zeta, \bigdt \Ad{\exp(tx)}(y) \bigg\rangle  \\
        & = \bigg\langle \bigdt \Ad{\exp(-tx)}(\zeta), y \bigg \rangle =\langle [-x,\zeta], y\rangle = \langle [\zeta,x], y \rangle.
    \end{align*}
    Moreover, \(\mathfrak{p}_\zeta=\mathfrak{l}_\zeta \oplus \mathfrak{u}_\zeta\) where \(\mathfrak{u}_\zeta\) is the nilpotent radical obtained from the Levi factor \(\mathfrak{l}_\zeta\) (cf. section \ref{section_parabolic_subalgebras} and proposition \ref{prop_charact_parabolic_subgroups}). Hence, by definition of \(\mathfrak{l}_\zeta\), if \(x \in \mathfrak{l}_\zeta\), \(\chi_\zeta([x,y])=\langle [\zeta,x], y \rangle =0\). By the direct sum decomposition above, it suffices to see that \(\chi_\zeta([x,y])\) for \(x \in \mathfrak{u}_\zeta\). Now, since \(\mathfrak{u}_\zeta\) is an ideal of \(\mathfrak{p}_\zeta\), we have \([\mathfrak{u}_\zeta, \mathfrak{p}_\zeta] \subset \mathfrak{u}_\zeta\), hence \([x,y]\in \mathfrak{u}_\zeta\) for \(x\in \mathfrak{u}_\zeta, y \in \mathfrak{p}_\zeta\). Therefore, we need only show that \(\chi_\zeta(x)=0\) for all \(x \in \mathfrak{u}_\zeta\), but this is trivially true because \(\mathfrak{u}_\zeta\) is the radical of the restriction of \(\langle \cdot, \cdot  \rangle\) to \(\mathfrak{p}_\zeta\) (that is, \(\langle \xi, x\rangle=0\) for all \(\xi \in \mathfrak{p}_\zeta, x \in \mathfrak{u}_\zeta\), see \parencite[Page 273]{knapp2013lie}).
\end{proof}

Consider now any parabolic subalgebra \(\mathfrak{p}\), and \(\chi\in \mathfrak{p}^*\) a character. The non-degenerate pairing \(\langle \cdot, \cdot \rangle\) induces a linear isomorphism \(\mathfrak{g}\cong \mathfrak{g}^*, \zeta \mapsto \langle \zeta, \cdot \rangle\). Recall that we have a direct sum decomposition \(\mathfrak{g}=\mathfrak{k}\oplus \ii \mathfrak{k}\), which induces a projection \(\mathfrak{g} \twoheadrightarrow \ii \mathfrak{k}\). Finally, the pairing \(\langle \cdot , \cdot \rangle\) induces a direct sum decomposition \(\mathfrak{g}=\mathfrak{p}\oplus \mathfrak{p}^\perp\). This comes with a natural orthogonal projection \(p_\perp \colon \mathfrak{g}\to \mathfrak{p}\), whose dual map gives an injection \(\mathfrak{p}^* \hookrightarrow \mathfrak{g}^*, \chi \mapsto \chi \circ \textsf{pr}_\perp \). We thus have the following chain of maps 
\[
\mathfrak{p}^* \hookrightarrow \mathfrak{g}^* \xrightarrow{\cong} \mathfrak{g} \to \ii \mathfrak{k}.
\]
Given \(\chi \in \mathfrak{p}^*\) we denote by \(\zeta_\chi\in \ii \mathfrak{k}\) its image through the previous chain of maps.

\Definition{}{Let \(\mathfrak{p}\subset \mathfrak{g}\) be a parabolic subalgebra. A character \(\chi \in \mathfrak{p}^*\) is said to be \textit{antidominant}\index{Character!antidominant} if \(\mathfrak{p}_{\zeta_\chi}\subset \mathfrak{p}\) and \textit{strictly antidominant}\index{Character!stricly antidominant} if \(\mathfrak{p}_{\zeta_\chi}=\mathfrak{p}\).}

\begin{obs} \label{obs_chis_antidominant}
    For each \(\zeta\in \ii \mathfrak{k}\), we have \(\zeta_{\chi_\zeta}=\zeta\), thus the character \(\chi_\zeta\) of \(\mathfrak{p}_\zeta\) is always strictly antidominant. This follows from the fact that \(\chi_\zeta \circ \textsf{pr}_\perp=\chi_\zeta\) (recall that \(\zeta \in \mathfrak{p}_\zeta\)) the way \(\langle \cdot, \cdot \rangle\) induces the pairing. The fact that each \(\chi_\zeta\) is antidominant is a very technical requirement, and we will not make explicit use of it in our sketch of the proof (but it is necessary nonetheless).
\end{obs}

Consider \(\mathcal{P}\) a holomorphic principal \(G\)-bundle over a Riemann surface \(X\). Let \(\zeta \in \ii \mathfrak{k}\) and let \(P_\zeta\subset G\) be the parabolic subgroup defined by \(\zeta\) as in proposition \ref{prop_charact_parabolic_subgroups}. Consider a reduction of the structure group of \(\mathcal{P}\) to \(P_\zeta\) given by a holomorphic section \(\sigma \in \mathcal{O}(X; \mathcal{P}/P_\zeta)\), and let \(\mathcal{P}_\sigma:=\sigma^*\mathcal{P}\) be the corresponding holomorphic \(P_\zeta\)-principal bundle over \(X\). Consider \(L_\zeta\) the Levi subgroup of \(P_\zeta\) defined by \ref{prop_charact_parabolic_subgroups}. Define \(K_\zeta:=K\cap L_\zeta\) which is a Lie group whose Lie algebra is given by \(\mathfrak{k}_\zeta:=\mathfrak{k}\cap \mathfrak{l}_\zeta\).  Then, \(K_\zeta\) is a maximal compact subgroup of \(L_\zeta\) and the inclusion \(K_\zeta \subset L_\zeta\) is a homotopy equivalence. The inclusion \(L_\zeta\subset P_\zeta\) is a homotopy equivalence as well, so \(K_\zeta\) and \(P_\zeta\) are homotopy equivalent. This implies that any reduction \(\sigma\) induces another reduction (unique up to homotopy equivalence) of the structure group to \(K_\zeta\). Call \(P'_\sigma\) the resulting smooth principal \(K_\zeta\)-bundle. See \parencite[Example 5.6]{kobayashi_vol1} and the results surrounding it for a justification of this fact.

\Definition{}{Let \(\mathcal{P}(X,G)\) be a holomorphic principal bundle and \(\zeta \in \ii \mathfrak{k}\). Keeping the notations of the previous paragraph, consider a connection \(A\) on \(\mathcal{P}'_\sigma\), and let \(F_A\in \Omega^2_\R(X; P'_\sigma\times_{K_\zeta} \mathfrak{k}_\zeta)\) be its curvature. We define the \textit{degree of \(\mathcal{P}\) with respect to \(\zeta\) and the reduction \(\sigma\)} by \label{def_degree_principal_bundle}
\begin{equation} \nonumber
    \deg \mathcal{P}(\sigma, \zeta)=\frac{\mathrm{i}}{2 \pi}\cdot \int_X \chi_\zeta(F_A).     
\end{equation}
}
\begin{obs} \label{obs_deg_well_defined}
    Let us check that \(\deg \mathcal{P}(\sigma, \zeta)\) is a well-defined real number. First, observe that \(\chi_\zeta(\mathfrak{k}_\zeta) \subset \ii \R\). Indeed, \(\zeta=\ii \zeta'\) for some \(\zeta' \in \mathfrak{k}\), hence \(\chi_\zeta(\xi)=\langle \ii \zeta', \xi \rangle = \ii \langle \zeta', \xi \rangle\) for \(\zeta \in \mathfrak{k}_\zeta \subset \mathfrak{k}\) using \(\C\)-bilinearity. Moreover, since \(\mathfrak{g}=\mathfrak{k}\oplus \ii \mathfrak{k}\), \(\mathfrak{g}\) is endowed with a natural conjugation map. With respect to this conjugation map, and due to \(\C\)-bilinearity, a direct computation shows that \(\overline{\langle \zeta, \xi \rangle} = \langle \overline{\zeta}, \overline{\xi} \rangle\) for any pair \(\zeta, \xi \in \mathfrak{g}\). Hence, if \(\zeta, \xi \in \mathfrak{k}\), they coincide with their complex conjugates and their product is real. 

    On the other hand, \(F_A\) takes values in the associated bundle \(P'_\sigma \times _{K_\zeta}\mathfrak{k}_\zeta\). We define the composition of \(\chi_\zeta\) and any \(p\)-form in \(\Omega^p_\R(X; P'_\sigma \times_{K_\zeta} \mathfrak{k}_\zeta)\) by linear extension of \(\chi_\zeta(\alpha \otimes \bm{\xi})=\chi_\zeta(\bm{\xi})\cdot\alpha\), where \(\bm{\xi}\in \Gamma(X; P'_\sigma \times_{K_\zeta}\mathfrak{k}_\zeta)\). Finally, if \(\bm{\xi}(x)=[p, \xi]\in P'_\sigma \times_{K_\zeta}\mathfrak{k}_\zeta\), with \(\pi(p)=x\), we define \(\chi_\zeta(\bm{\xi}(x))=\chi_\zeta(\xi)\). This does not depend on the representative chosen. Indeed, \([p, \xi]=[p\cdot k, \Ad{k^{-1}}(\xi)]\) for \(k \in K_\zeta\), but
    \[
    \chi_\zeta(\Ad{k^{-1}}(\xi))=\langle \zeta, \Ad{k^{-1}}(\xi) \rangle = \langle \Ad{k}(\zeta), \xi \rangle = \langle \zeta, \xi \rangle = \chi_\zeta(\xi)
    \]
    using the \(\Ad{}\)-invariance of the product and the fact that \(K_\zeta \subset L_\zeta=\corch{g\in G: \Ad{g}(\zeta)=\zeta}\). Finally, let us see that the definition is independent of the chosen connection \(A\). First, recall that since the space of connections \(\mathscr{A}(P'_\sigma)\) is an affine space modelled on \(\Omega^1_\R(X; \mathrm{ad}(P_\sigma'))\), if we fix a connection \(A \in \mathscr{A}(P_\sigma')\), any other connection will be given by \(A+\alpha\) for some \(\alpha \in \Omega_\R^1(X; P_\sigma')\). Observe that by our definitions, and due to the fact that everything is smooth, it follows that \(\chi_\zeta(\alpha)\in \Omega_\R^1(X)\). Let us assume that \(\chi_\zeta(\diff_A \alpha)=\diff \chi_\zeta(\alpha)\) (we give a proof of this fact in lemma \ref{lemma_def_degree_principal_bundl} below). The curvature two-forms of \(A\) and \(A+\alpha\) relate as \(F_{A+\alpha}=F_A+\diff_A \alpha + 1/2 \cdot [\alpha \wedge \alpha]\), (see lemma \ref{lemma_F_A_plus_alpha}). Since \(\chi_\zeta\) is a character of \(\mathfrak{p}_\zeta\), it vanishes on \([\mathfrak{p}_\zeta, \mathfrak{p}_\zeta]\), hence, it vanishes on \([\mathfrak{k}_\zeta, \mathfrak{k}_\zeta]\) as well, and that is where \([\alpha \wedge \alpha]\) takes values. Hence, applying Stokes' theorem, we find
    \[
    \int_X \chi_\zeta(F_{A+\alpha})=\int_X \chi_\zeta(F_A) + \int_X \diff \chi_\zeta(\alpha)
    \]
    the latter being zero because \(\partial X=\emptyset\).
\end{obs}

\Lemma{}{Under the conditions of definition \ref{def_degree_principal_bundle}, we have \(\chi_\zeta(\diff_A \alpha)=\diff \chi_\zeta(\alpha)\) for all \(\zeta\in \ii \mathfrak{k}\) and \(\alpha \in \Omega_\R^1(X; \mathrm{ad}(P_\sigma'))\). \label{lemma_def_degree_principal_bundl}}
\begin{proof}
    The calculation is similar to that of lemma \ref{lemma_F_A_plus_alpha}. Let \(\pi\colon P_\sigma'\to X\) be the projection, and \(\alpha \in \Omega_\R^1(X; P_\sigma' \times_{K_\zeta} \mathfrak{k}_\zeta)\). Recall from theorem \ref{th_forms_associated_bundles} that \(\Upsilon^{-1}(\alpha) \in \Omega^1_{\text{Ad, hor}}(P'_\sigma; \mathfrak{k}_\zeta)\). Expand \(\alpha=\alpha^j \otimes \bm{\xi}_j\) for some \(\alpha^j \in \Omega_\R^1(X)\) and \(\bm{\xi}_j\in \Gamma(X; P_\sigma'(\mathfrak{k}_\zeta))\). As in lemma \ref{lemma_F_A_plus_alpha} \(\Upsilon^{-1}(\alpha)=\pi^* \alpha^j \otimes \hat{\xi}_j\), where \(\hat{\xi}_j=\Upsilon^{-1}(\bm{\xi}_j)\colon P'_\sigma \to \mathfrak{k}_\zeta\) is \(K_\zeta\)-equivariant. After setting this notation up, the result follows from a simple computation. Let \(p \in P_\sigma'\) and \(v_1, v_2 \in T_pP_\sigma'\). On the one hand, it follows from our definitions that \(\diff \chi_\zeta(\alpha)=\chi_\zeta(\bm{\xi}_j)\cdot\diff \alpha^j\). Hence, 
    \begin{align*}
        \diff_{\pi(p)}\ \chi_\zeta(\alpha) (\diff_p \pi(v_1), \diff_p\pi(v_2)) & = \diff_{\pi(p)} \alpha(\diff_p\pi(v_1), \diff_p\pi(v_2)) \cdot \chi_\zeta(\bm{\xi}_j(\pi(p)))\\
        &=(\pi^*\alpha^j)_p(v_1, v_2) \cdot \chi_\zeta(\bm{\xi}_j(\pi(p))).
    \end{align*} 
    On the other hand,
    \begin{align*}
        \chi_\zeta((\diff_A \alpha)_p(\diff_p\pi(v_1), \diff_p\pi(v_2))) & = \chi_\zeta(q_p( \Diff_A \Upsilon^{-1}(\alpha))_p(v_1^H, v_2^H)) = \chi_\zeta\big(q_p(\diff_p(\pi^* \alpha^j)(v_1^H,v_2^H)\cdot \hat{\xi}_j(p))\big) \\
        & = (\pi^* \diff\alpha^j)_p(v_1, v_2) \cdot \chi_\zeta(\bm{\xi}_j(\pi(p))).
    \end{align*}
    In the last step we used the fact that \(\xi_j(\pi(p))=[p, \hat{\xi}_j(p)]\) and the fact that \(\diff_p\pi(v)=\diff_p\pi(v^H)\) because vertical vectors are annihilated by the projection. We thus get the identity \(\diff \chi_\zeta(\alpha)=\chi_\zeta(\diff_A \alpha)\) but written pointwise.
\end{proof}

We are now ready to state the definition of stability. First, observe that parabolic subgroups \(P_\zeta\) are embedded subgroups, and the action of \(P_\zeta\) on \(G\) by left multiplication \(P_\zeta \times G \to G, (p,g) \mapsto p\cdot g \) is holomorphic free and proper (the reasoning in the proof of \parencite[Th. 21.17]{lee_introSM} holds), hence, by theorem \ref{quot_man_th_complex}, \(G/P_\zeta\) has the structure of a complex manifold. Then, given a holomorphic principal bundle \(\mathcal{P}\), the associated bundle \(\mathcal{P}(G/P_\zeta)\) is holomorphic as well, and we can consider its holomorphic sections.

\Definition{\textbf{Stability condition}}{Let \(\mathcal{P}_G\) be a holomorphic principal bundle over a Riemann surface, and let \(c \in \ii \mathfrak{z}_K\). We say that \(\mathcal{P}_G\) is \label{def_stability_connected}
\begin{enumerate}
    \item \(c\)\textit{-stable} (resp. \(c\)\textit{-semistable}) if for any \(\zeta \in \ii \mathfrak{k}\), and any \(\sigma \in \mathcal{O}(X; \mathcal{P}(G/P_\zeta))\), we have 
    \[
    \deg \mathcal{P}(\sigma, \zeta)-\langle c, \zeta\rangle >0 \hspace{1.5 cm} (\text{resp. } \deg \mathcal{P}(\sigma, \zeta)-\langle c, \zeta \rangle \geq 0).
    \]
    \item \(c\)\textit{-polystable} if it is \(c\)-semistable and or any \(\zeta\in \ii \mathfrak{k}\) and any \(\sigma \in \mathcal{O}(X; \mathcal{P}(G/P_\zeta))\) such that \(\deg \mathcal{P}(\sigma, \zeta)-\langle c, \zeta \rangle = 0\), there is a holomorphic reduction of the structure group \(\sigma'\in \mathcal{O}(X; \mathcal{P}_\sigma(P/L_\zeta))\).
\end{enumerate}
}

The arguments of remark \ref{obs_deg_well_defined} apply to show that \(\langle c, \zeta \rangle \in \R\), hence the inequalities above make sense. With the correct notion of stability, we can now state the correspondence.

\begin{obs}
    If \(X\) is not a Riemann surface but a compact Kähler manifold, the definition is slightly modified because the sections one has to consider are not defined on all of \(X\), but on open sets \(X_0\) such that \(X\setminus X_0\) is a complex submanifold satisfying \(\operatorname{codim}_\C( X\setminus X_0)\geq 2\). 
\end{obs}

\Theorem{\textbf{Hitchin-Kobayashi for principal bundles}}{Let \(G\) be a connected complex reductive Lie group, \(K\subset G\) a maximal compact subgroup, and \(\mathcal{P}_G\) a holomorphic principal \(G\)-bundle over a compact Riemann surface \(X\). Then, \(\mathcal{P}_G\) is \(c\)-polystable if and only if there is a reduction of the structure group to \(K\), \(P_K\) whose Chern connection \(A\) satisfies the equation  \label{th_HK_principal_bund_connected}
\begin{equation} \label{eq_Hermite_Yang_Mills}
    \Lambda_\omega F_A=-\ii  c
\end{equation}
for some \(c \in\ii\mathfrak{z}_K\). 
}
Equation \eqref{eq_Hermite_Yang_Mills} is called the \textit{hermitian Yang-Mills equation for principal bundles}. We insist on the fact that \(F_A \in \Omega^2_\R(X; P_K \times_K \mathfrak{k})\) is the curvature of the Chern connection on the \textit{reduced} bundle.

\subsection{Infinite-dimensional analogue of the Kempf-Ness theorem} \label{section_inf_dim_Kempf_Ness}
We have claimed several times that the Hitchin-Kobayashi correspondence can be seen as an infinite-dimensional version of the Kempf-Ness theorem. Section \ref{section_moment_map_picture} was devoted to establish, with some level of detail, the moment map picure for the moduli space of solutions of the hermitian Yang-Mills equation. In this section, we shall explore how this situation can be seen as a formal analogue of the Kempf-Ness theorem when the base manifold is a Riemann surface \(X\). We will not try to give proofs for most of the claims stated in this section because their proofs require very advanced techniques coming from the fields of analysis and elliptic differential equations, and we have not developed the tools required to deal with them. There are two points we should clarify. First, at the date of writing of this document \textit{there is no infinite-dimensional version of the Kempf-Ness theorem}. Hence, by saying that the Hitchin-Kobayashi correspondence is an infinite-dimensional analogue of said theorem, we mean that it is a proof of the statement of the Kempf-Ness theorem for the very specific case of the action of the gauge group on the space of connections. The second point is that this theorem is a \textit{formal analogue}: it should not be interpreted as a word-for-word version of the finite-dimensional version, but rather as a philosophy, as we shall discuss below. Along the lines of this second point, there are a few things we should remark. The first one is that when we say that \(\mathscr{G}^\C=\Gamma(X; \mathrm{ad}(P_G))\) is a complexification of \(\mathscr{G}=\Gamma(X; \mathrm{Ad}(P_K))\) we do not mean it literally in the sense of satisfying a universal property as in definition \ref{def_complex_Lie_group}. We can \textit{formally} say that it is a complexification because fiberwise, it is. Secondly, if one looks at the definition of stability for principal bundles we have given here, and the definition we gave for the finite-dimensional version of the Kempf-Ness theorem (cf. definitions \ref{def_stability_connected} and \ref{def_stability_fin_dim_KempfNess}) it does not look at all like one is a generalization of the other. There are ways to write these conditions that make them look more similar to each other, but that does not matter much to see how \textit{the philosophy} of the theorem is preserved. For the rest of the section, we fix a principal \(K\)-bundle, and set \(P_G:=P_K \times_K G\), which is a principal \(G\)-bundle. 

For the algebraic side, recall that the Chern-Singer correspondence induces a bijection between the set of holomorphic structures on \(P_G\), \(\mathscr{J}(P_G)\), and the space of connections on \(P_K\), \(\mathscr{A}(P_K)\). Polystability is a condition depending on the holomorphic structure, but by means of the Chern-Singer correspondence, we may define a connection \(A \in \mathscr{A}(P_K)\) as polystable if \(\mathsf{CS}^{-1}(A)=J\) makes \(P_G\) into a polystable holomorphic principal bundle. Hence, we can define \(\mathscr{A}^{ps}(P_K) \subset \mathscr{A}(P_K)\). Since the action of \(\mathscr{G}^\C\) on \(\mathscr{A}(P_K)\) is induced by the action of \(\mathscr{G}^\C\) on \(\mathscr{J}(P_G)\), and polystability is preserved under biholomorphism, we deduce that \(\mathscr{G}^\C\) preserves \(\mathscr{A}^{ps}(P_K)\). Hence, we may consider \(\mathscr{A}^{ps}(P_K)/\mathscr{G}^\C\). This quotient then represents the set of holomorphic polystable structures on \(P_G\) up to biholomorphism. 

Now, the Hitchin-Kobayashi correspondence can be interpreted as defining a bijection
\begin{equation} \label{eq_bijec_HK}
\mu^{-1}(0)/\mathscr{G}_{k+1} \longleftrightarrow \mathscr{A}^{ps}(P_K)/\mathscr{G}^\C,
\end{equation}
and this very clearly mimics the Kempf-Ness theorem: we have a Hamiltonian action of a group on a manifold, and there is a bijection between the symplectic quotient and the quotient of a set of polystable points by the complexification of the gauge group. Observe that the Kempf-Ness theorem can be interpreted as saying that the symplectic quotient is homeomorphic to the usual topological quotient we get from the action of the complexification of the group when the set of ``bad points'' (i.e. non-polystable points) is removed. Here we have the same thing: there is a bijection between the symplectic quotient and the set of connections quotiented by the complexified gauge group where the set of ``bad points'', that is, non-polystable points, has been removed.

It remains to see how theorem \ref{th_HK_principal_bund_connected} defines the bijection in \eqref{eq_bijec_HK}. Given an equivalence class of connections in \(\mu^{-1}(0)/\mathscr{G}\), choose a smooth connection \(A\). This connection defines an integrable almost holomorphic structure on \(P_G\) through the Chern-Singer correspondence, and \(A\) is clearly the Chern connection with respect to that almost holomorphic structure. Since \(A\) satisfies equation \eqref{eq_Hermite_Yang_Mills}, the Hitchin-Kobayashi correspondence implies that the almost holomorphic structure it endows is polystable. Any other smooth connection \(A'\) we may find in the same orbit is related to \(A\) by \(A'=\phi \cdot A\) for some \(\phi \in \mathscr{G}\). Since \(\mathscr{G}\) is a subgroup of \(\mathscr{G}^\C\), and the action of the latter on the space of connections extends the action of \(\mathscr{G}\), we get two integrable almost holomorphic structures in the same \(\mathscr{G}^\C\), hence, the map is well-defined. Observe that we are basically sending \(\mathscr{G}\)-orbits to \(\mathscr{G}^\C\)-orbits. 

To see that the construction above actually defines a bijection, we need to shift our viewpoint in the Hitchin-Kobayashi correspondence. The statement we have given in theorem \ref{th_HK_principal_bund_connected} fixes a holomorphic structure on \(P_G\) and looks for reductions of the structure group to \(K\) whose Chern connection satisfies equation \eqref{eq_Hermite_Yang_Mills}. However, the bijection given by \eqref{eq_bijec_HK} is phrased differently. There, we \textit{are fixing a reduction of the structure group}, \(P_K\), and we \textit{are looking for polystable holomorphic structures} on \(P_G\), or equivalently, connections \(A\in \mathscr{A}(P_K)\) satisfying equation \eqref{eq_Hermite_Yang_Mills}. Theorem \eqref{th_HK_principal_bund_connected} can be stated as follows (this is similar to theorem \ref{th_HK_simple_connections} below): Fix \(P_K\) a principal \(K\)-bundle, and \(P_G=P_K \times_K G\) a principal \(G\)-bundle. Let \(A \in \mathscr{A}(P_K)\) be a connection. Then, there is a gauge transformation \(\phi \in \mathscr{G}^\C\) such that \(\phi\cdot A\) satisfies \eqref{eq_Hermite_Yang_Mills} if and only if \(P_G\) is polystable with respect to the holomorphic structure given by \(\mathsf{CS}^{-1}(A)\). 

We can now show that \eqref{eq_bijec_HK} is a bijection. Given \(A \in \mathscr{A}^{ps}(P_K)\), by the formulation above, there is \(\phi \in \mathscr{G}^\C\) such that \(A':=\phi\cdot A\) satisfies equation \eqref{eq_Hermite_Yang_Mills}, hence \(A'\in \mu^{-1}(0)\). Then the equivalence class of \(A'\) is mapped to that of \(A\). For injectivity, let \(A', A''\in \mu^{-1}(0)\) be in the same \(\mathscr{G}^\C\)-orbit, thus, there is \(\psi\in \mathscr{G}^\C\) such that \(A'=\psi\cdot A''\), but one may actually show that \(\psi\in \mathscr{G}\), hence they are in the same \(\mathscr{G}\)-orbit. Again, this is along the lines of theorem \ref{th_HK_simple_connections} below.

When we consider a compact Kähler manifold instead of a Riemann surface, we look for connections satisfying \(\Lambda_\omega F_A=-\ii c\) \textit{and} \(F_A^{(0,2)}=0\) (so that they induce a holomorphic structure on \(P_G\) through the Chern correspondence). When constructing the moduli space of solutions of the hermitian Yang-Mills equations there is a very important difference from what we did in chapter \ref{chapter_HYM_equation}: the space \(\mathscr{A}^{(1,1)}\) is not an affine space. In fact, it is not even a manifold because it may have singular points. It is a Banach manifold outside its singularity set, and there \(\mathscr{G}\) acts symplectically for the restriction of the moment map. The moduli space of solutions can still be constructed using Banach-symplectic reduction.

\subsection{A sketch of the proof} \label{section_sketch_proof}
We will sketch a proof of theorem \ref{th_HK_principal_bund_connected} following Mundet i Riera, García Prada and Gothen \parencite{i1999hitchin} and \parencite{garcia2009hitchin}. The results stated in these articles are more general: they are given in the context of Kähler fibrations and Higgs bundles respectively, but we may reduce their arguments to ours by ignoring the extra structures. For an alternative proof of this result which relies heavily on the vector bundle theorem, we refer the reader to \parencite{anchouche2001einstein}. Our strategy for the proof is to reduce the structure group further to another subgroup in such a way that another stronger condition is satisfied: simplicity. Then, the results will follow from the arguments in \parencite{i1999hitchin}. The results in these articles are quite involved; giving detailed proofs would take a lot of work, and it is beyond the scope of this work. We include a sketch for the sake of completeness.  

First of all, we say that an element \(\bm{\xi} \in \mathrm{Lie}(\mathscr{G}^\C)=\Gamma(X; P_G \times_G \mathfrak{g})\) is \textit{semisimple} if, for any \(x \in X\), \(\bm{\xi}(x)\in \mathfrak{g}\) is semisimple (i.e. its adjoint diagonalizes). We see \(\bm{\xi}(x)\in \mathfrak{g}\) through any of the isomorphisms \((P_G \times_G \mathfrak{g})_x \cong \mathfrak{g}\). The condition of being semisimple does not depend on the chosen representative (because any other representative will be its image through the adjoint action, and \(\Ad{g}\) is an isomorphism for every \(g \in G\)). Finally, recall that the action of \(\mathscr{G}\) on \(\mathscr{A}\) can be extended to an action of \(\mathscr{G}^\C\) using the Chern-Singer correspondence (see section \ref{section_Chern_Singer_correspondence}). Then, a connection \(A \in \mathscr{A}(P_G)\) is said to be \textit{simple} if no element \(\bm{\xi} \in \mathrm{Lie}(\mathscr{G})\) fixes \(A\) under the infinitesimal action of \(\mathscr{G}^\C\). More explicitly, given \(\bm{\xi}\in \mathrm{Lie}(\mathscr{G}^\C)\) we can construct the associated vector field \(X_{\bm{\xi}}\). Then, the \textit{infinitesimal action} is the vector \(X_{\bm{\xi},A}\in \Omega^1(M;\mathrm{ad}(P_G))\). As shown in \parencite[Prop. 2.3]{trautwein2015surveygitpictureyangmills}, this is given, explicitly, by
\[
X_{\bm{\zeta}_1+ \ii \bm{\zeta}_2, A}=-\diff_A \zeta_1 - \star \  \diff_A \zeta_2
\]
for \(\bm{\xi}=\bm{\zeta}_1 + \ii \bm{\zeta}_2\) and \(\bm{\zeta}_1 , \bm{\zeta}_2 \in \mathrm{Lie}(\mathscr{G})\). Then, \parencite{i1999hitchin} shows the following version of the Hitchin-Kobayashi correspondence. By ``not fixing \(A\) under the infinitesimal action'' we mean \(X_{\bm{\xi}, A} \neq 0\).

\Theorem{\textbf{Hitchin-Kobayashi for simple connections}}{ Let \(P_K\) be a principal \(K\)-bundle over a Riemann surface, and \(P_G :=P_K \times_K G\). Let \(A \in \mathscr{A}(P_K)\) be a simple connection and \(c \in \ii \mathfrak{z}_K\). Then, there is a gauge transformation \(\psi \in \mathscr{G}^\C\) such that \label{th_HK_simple_connections}
\begin{equation} \label{eq_HK_simple}
\Lambda_\omega F_{\psi \cdot A}=-\ii c
\end{equation}
if and only if \(P_G\) is \(c\)-stable with respect to the integrable almost holomorphic structure \(\mathsf{CS}^{-1}(A)\). If \(\psi\cdot A\) and \(\psi'\cdot A\) solve equation \eqref{eq_HK_simple}, then there is some \(\phi \in \mathscr{G}\) such that \(\psi'=\phi \circ \psi\).
}
\begin{proof}[\textit{Sketch of proof } \parencite{i1999hitchin}]
    The main idea is to construct a functional \(\Psi\colon \mathscr{A}_k(P_K) \times \mathscr{G}_{k+1}^\C \to \R\) whose critical points are solutions of \eqref{eq_HK_simple} after completing these spaces with suitable Sobolev norms as in section \ref{section_moment_map_picture}. We fix a \textit{hermitian metric} \(\langle \cdot, \cdot \rangle\) in \(\mathfrak{g}\). This is done by means of a representation as follows. First, since \(K\) is compact, it admits a faithful representation \(\rho\colon K \to \mathrm{U}(V)\) where \(V\) is some hermitian vector space. By the universal property of complexifications, this extends to a faithful representation \(\rho_\C\colon G \to \mathrm{GL}(V)\). We may then define a hermitian inner product on \(\mathfrak{g}\) using this representation as follows
    \[
    \langle \xi, \zeta \rangle:=\mathrm{Tr}(\diff_e\rho_\C(\xi)\cdot \diff_e\rho_C(\zeta)^*).
    \]
    This extends to a hermitian fiber metric on \(P_G \times_G \mathfrak{g}\). 

    First, one defines the \textit{weight} (\((\cdot, \cdot)\) is the pairing between \(\mathrm{Lie}(\mathscr{G}_{k+1})\) and its dual)
    \[
    \lambda_t(A; \bm{\zeta}):=\mu^{\bm{\zeta}}(\exp_{\mathscr{G}_{k+1}}(t\  \bm{\zeta})\cdot A)=\big(\mu(\exp_{\mathscr{G}_{k+1}}(t \ \bm{\zeta})\cdot A) , \  \bm{\zeta} \big) \in \R, \hspace{1 cm} \bm{\zeta} \in \mathrm{Lie}(\mathscr{G}_{k+1}).
    \]
    One may show that \(\lambda_t\) increases with \(t\in \R\), hence, we define (and it always exists) the \textit{maximal weight}
    \[
    \lambda(A; \bm{\zeta})=\lim_{t \to 0} \lambda_t(A; \bm{\zeta}) \in \R \cup \corch\infty.
    \]
    Then, one shows that there is a unique functional \(\Psi\colon \mathscr{A}_k(P_K) \times \mathscr{G}_{k+1}^\C \to \R\), called \textit{the integral of the moment map} satisfying:
    \begin{enumerate}
        \item For any \(A \in \mathscr{A}_k(P_K)\), the critical points of the restriction \(\Psi_A\colon \mathscr{G}_{k+1} \to \R\) of \(\Psi\) to \(\corch{A} \times \mathscr{G}_{k+1}\) coincide with the points of the orbit \(\mathscr{G}_{k+1}\cdot A\) on which \(\mu\) vanishes.
        \item The restriction of \(\Psi_A\) to lines of the form \(\corch{\exp_{\mathscr{G}_{k+1}}(t\bm{\zeta}): t \in \R}\) for \(\bm{\zeta} \in \mathrm{Lie}(\mathscr{G}_{k+1})\) is convex.
    \end{enumerate}
    One then shows that an element \(\psi \in \mathscr{G}_{k+1}\) is a critical point of \(\Psi_A\) for some \(A \in \mathscr{A}_k(P_K)\) if and only if \(\mu(\psi\cdot A)=0\), i.e. if \(\psi\cdot A\) is a solution of \eqref{eq_HK_simple}.  
    % Then, one shows that the integral of the moment map may be written as
    % \[
    % \Psi_A(\exp(i\bm{\zeta}))=\int_X \langle \Lambda_\omega F_A+\ii c, \bm{\zeta}\rangle +\int_0^1(1-\ell)\cdot ||\exp(i\ell \bm{\zeta})\ \barpartial_A(\bm{\zeta}) \ \exp(-i \ell \bm{\zeta})||^2 \ \diff \ell. 
    % \]
    % In the previous expression, \(||\cdot ||^2:=\int_X |\cdot |^2\), where \(|\cdot |^2\) is the norm induced by the hermitian product \(\langle \cdot, \cdot \rangle\) on \(\mathfrak{g}\). Moreover, \(\barpartial_A=\Pi^{(0,1)}\circ \diff_A\) is the \((0,1)\) component of the induced connection \(\diff_A\) on the associated bundle \(P_K \times_K \mathfrak{k}\), and \(|\cdot|^2\) is the norm associated to the \(\C\)-bilinear product \((\cdot, \cdot)\).  
    
    After everything is set up, we can show that stability implies the existence of solutions. Let \(A\) be a simple connection. One may show that \(\Psi(A, \phi \circ  \psi)=\Psi(A, \psi)\) for \(\phi \in \mathscr{G}_{k+1}\) and \(\psi \in \mathscr{G}^\C_{k+1}\) \parencite[Prop. 3.4]{i1999hitchin}. Then, \(\Psi_A\) factors through \(\mathscr{G}_{k+1}\):
    \[
    \Psi_A\colon \mathscr{G}^\C_{k+1}/\mathscr{G}_{k+1} \to \R.
    \]
    We want to minimize that functional. To do so, since \(\mathfrak{g}/\mathfrak{k}\cong \ii \mathfrak{k}\), we may use the exponential map to identify \(\mathscr{G}^\C/\mathscr{G}\) with \(\Gamma(M;P_K \times_K \ii \mathfrak{k})\). Choosing \(p>2\), we define \(\mathscr{M}^p_2=L_2^p(P_K \times_K \ii \mathfrak{k})\) (the \(L_2^p\) norm is defined as usual with respect to the hermitian fiber product defined as above), which we shall identify with the quotient of the Sobolev completions. When \(A\) is a simple connection, we can restrict ourselves to looking for minima on the subset
    \[
    \mathscr{M}^p_{2,B}:=\corch{\bm{\zeta}\in \mathscr{M}^p_2:||\mu(\exp(\bm{\zeta})\cdot A)||^p_{L_2^p} \leq B}
    \] 
    for some positive constant \(B\). We see the moment map as taking values in \(\mathrm{Lie}(\mathscr{G}_{k+1})\) to make sense of the norm (cf. remark \ref{obs_mu_lie_sin_dual}), that is 
    \[
    ||\mu(\exp(\bm{\zeta})\cdot A)||^p_{L_2^p}=||\Lambda_\omega F_{\exp_{\mathscr{G}_{k+1}}(\bm{\zeta})\cdot A}+\ii c||_{L_2^p}^p.
    \] 
    Finally, one shows that there is an element of \(\mathscr{Met}^p_{2,B}\) that minimizes \(\Psi_A\) and this minimizing element is smooth. Then, it corresponds to an equivalence class of the quotient \(\mathscr{G}^\C_{k+1}/\mathscr{G}_{k+1}\), so there is a solution, and any two different solutions are related by a transformation of \(\mathscr{G}_{k+1}\). This is rather technical, and it is the content of section \parencite[Sec. 6.2]{i1999hitchin}.   

    Conversely, choose a simple connection \(A\) and a gauge transformation \(\psi \in \mathscr{G}^\C_{k+1}\) such that \(\psi\cdot A\) satisfies \eqref{eq_HK_simple}. To show that \(A\) is stable, let \(\zeta \in \ii \mathfrak{k}\), and choose any holomorphic reduction of the structure group to \(P_\zeta\), \(\sigma \in \mathcal{O}(X; \mathcal{P}(G/P_\zeta))\). By definition, one should check that \(\deg \mathcal{P}(\sigma, \zeta)-\langle c, \zeta\rangle >0\), however, \parencite[Lemma 4.3]{i1999hitchin} shows that one may equivalently check that \(\lambda(A;-\ii g_{\sigma,\zeta})>0\), where \(g_{\sigma, \zeta} \in \ii \mathrm{Lie}(\mathscr{G})=\Gamma(X; P_K \times_K \ii \mathfrak{k})\) is fiberwise dual to the character \(\chi_\zeta\) with respect to the hermitian fiber product induced by \(\langle \cdot, \cdot \rangle\) (see \parencite[Sec. 2.7]{i1999hitchin}). We proceed by contradiction. First, one can show the following bound \parencite[Sec. 6.3]{i1999hitchin}
    \begin{equation} \label{eq_dem_HK_mundet}
    t \sup |g_{\sigma, \zeta}| \leq C_1 \cdot \Psi_A(\exp_{\mathscr{G}_{k+1}}(t g_{\sigma,\zeta})) + C_2, \hspace{1 cm} C_1, C_2>0,
    \end{equation}
    and for sufficiently large \(t>0\). Moreover, the following relation can be proven 
    \[
    \Psi_A(\exp_{\mathscr{G}_{k+1}}( t g_{\sigma, \zeta}))=\int_0^1 \lambda_{t\cdot\ell} (A; -\ii g_{\sigma, \zeta}) \ \diff \ell.
    \]
    If \(\lambda(A; -\ii g_{\sigma, \zeta})<0\), then the previous integral would be negative for every \(t\) (because \(\lambda_t\) increases with \(t\)). Finally, \eqref{eq_dem_HK_mundet} would imply \(t\cdot \sup|g_{\sigma, \zeta}|<C_2\), but the left-hand side tends to \(+\infty\), so we get a contradiction.  

    We refer to \parencite{i1999hitchin} for the uniqueness claim.
\end{proof}

Since our goal is to use theorem \ref{th_HK_simple_connections} to prove theorem \ref{th_HK_principal_bund_connected}, we will need a way to translate the stability condition into a condition dealing with simplicity. We do this by means of the following theorem, which associates to each \(c\)-polystable principal bundle a \(c\)-stable principal bundle over a different group.

\Theorem{\textbf{Jordan-Hölder reduction}}{Let \(\mathcal{P}_G\) be a holomorphic principal \(G\)-bundle over a compact Riemann surface, and let \(c\in \ii \mathfrak{z}_K\). Then, there is a compact subgroup \(K'\subset K\) and a holomorphic reduction of the structure group of \(\mathcal{P}_G\) to \(G':=(K')^\C\), that is \(c\)-stable. Any two such reductions \(\mathcal{P}^{(1)}_{G'}\) and \(\mathcal{P}_{G'}^{(2)} \subset \mathcal{P}_G\) are related by \(\mathcal{P}_{G'}^{(1)}=\mathcal{P}_{G''}^{(2)}\cdot g\) for some \(g \in G\).  \label{th_Jordan_Holder}}
\begin{proof}
    The proof is not trivial at all. We refer the reader to section 2.10 of \parencite{garcia2009hitchin}.
\end{proof}

\begin{proof}[\textit{Sketch of proof of theorem} \ref{th_HK_principal_bund_connected}]
    Suppose first that \(\mathcal{P}_G\) is \(c\)-stable, and fix a reduction of the structure group to \(K\), \(P_K\). Then, one can show that it admits a simple connection \(A \in \mathscr{A}(P_K)\). The holomorphic structure induced by \(A\) on \(P_G\) via the Chern-Singer correspondence coincides with that of \(\mathcal{P}_G\). By theorem \ref{th_HK_simple_connections}, there is \(\varphi \in \mathscr{G}^\C\) such that \(A':=\varphi\cdot A\) satisfies equation \ref{eq_Hermite_Yang_Mills}. Now, since \(A\) and \(A'\) are on the same \(\mathscr{G}^\C\)-orbit, they induce the same structure of holomorphic principal bundle on \(P_G\) (see \eqref{CS_corresp_quot}). Finally, \(A'\) is the Chern connection with respect to \(\varphi\cdot J\), \(J=\mathsf{CS}^{-1}(A)\), and \(\varphi\cdot J\) induces the same structure of holomorphic principal bundle on \(P_G\) as \(J\). 
    
    Suppose now that \(\mathcal{P}_G\) is \(c\)-polystable, but not stable. We use theorem \ref{th_Jordan_Holder} and obtain a \(c\)-stable reduction of the structure group to some subgroup \(G'\). Now, \(\mathcal{P}_{G'}\) is stable, and the previous reasoning applies. Hence, there is a reduction of the structure group to \(P_{K'}\) whose Chern connection satisfies \eqref{eq_HK_simple}. Now we define \(P_K:=P_{K'} \times _{K'} K\subset P_{K'} \times_{K'} G=P_G\) (cf. \ref{obs_red_struc_group}) for the action of \(K'\) on \(K\) by inverse left multiplication. This is a reduction of the structure group to \(K\), and one can show that its Chern connection satisfies \eqref{eq_Hermite_Yang_Mills}. 

    For the converse, one shows first that if \(\mathcal{P}_G\) admits a reduction of the structure group to \(K\) satisfying \eqref{eq_Hermite_Yang_Mills}, then it is \(c\)-semistable (see \parencite{bradlow2003relative}). Let \(P_K\) be the corresponding \(K\)-bundle. Now, to prove polystability, let \(\sigma \in \Gamma(X; P(G/K))\) be a reduction of the structure group satisfying \eqref{eq_Hermite_Yang_Mills}. We know that the bundle is semistable, thus, for some parabolic subgroup \(P_\zeta \subset G\), we consider a holomorphic reduction of the structure group \(\sigma \in \mathcal{O}(X, \mathcal{P}(G/P_\zeta))\) such that \(\deg \mathcal{P}(\sigma, \zeta)=\langle c, \zeta \rangle\). We seek a further restriction of the structure group to the Levi subgroup \(L_\zeta\). Let \(P_K\) be the reduction of the structure group whose Chern connection satisfies \eqref{eq_Hermite_Yang_Mills}. This reduction produces a section \(\bm{\zeta}\) of the bundle \(P_K \times_K \ii \mathfrak{k}\). Since \(\mathrm{Lie}(\mathscr{G})=\Gamma(X;P_K \times_K \mathfrak{k})\), we can identify \(-\ii \bm{\zeta}\) with an element of \(\mathrm{Lie}(\mathscr{G})\). Now, let \(J\) be the almost complex structure of \(\mathcal{P}_G\), and let \(A=\mathsf{CS}(J)\in \mathscr{A}(P_K)\). One can see that the condition \(\deg \mathcal{P}(\sigma, \zeta)=\langle c, \zeta\rangle\) is equivalent to 
    \[
    \lambda(A, -\ii \bm{\zeta}):=\lim_{t \to \infty} \lambda_t(A, -\ii \bm{\zeta})=\lim_{t\to \infty} ( \mu(\exp_{\mathscr{G}}(t \bm{\zeta})\cdot A), -\ii \bm{\zeta}) =0.
    \]
    From \parencite[Lemma 2.2]{i1999hitchin}, we know that the function \(\lambda_t(A, -\ii \bm{\zeta})\) is non-decreasing. Moreover, by the way the moment map \(\mu\) is defined, equation \eqref{eq_Hermite_Yang_Mills} is equivalent to the vanishing of the moment map. Hence, at \(t=0\), we have \(\mu(\exp_{\mathscr{G}}(t \bm{\zeta})\cdot A)|_{t=0}=\mu(A)\), and the latter vanishes because the Chern connection \(A\) satisfies \eqref{eq_Hermite_Yang_Mills}. Therefore, we have \(\lambda_0(A, -\ii \bm{\zeta})=0\). Now, the fact that the weight is non-decreasing implies that it is constant and equal to \(0\) for any \(t \geq 0\). One may deduce from this fact that \(\exp_{\mathscr{G}}(t \bm{\zeta})\cdot J=J\) for all \(t \geq 0\). This implies that \(\bm{\zeta}\) induces a holomorphic section \(\sigma_\zeta \in \mathcal{O}(M; \mathcal{P}_\sigma(P_\zeta/L_\zeta))\), hence a reduction of the structure group to \(L_\zeta\), and this is precisely the definition of polystability.
\end{proof}

\section{The HK correspondence for non-connected groups}
This section contains the most original part of the document. It formulates a Hitchin-Kobayashi correspondence for principal bundles with non-connected complex reductive structure groups. Most of the section is based on the unpublished article \parencite{HKnonconnected} that my supervisor kindly shared with me. His work is more general, but with his ideas, we try to outline a generalization of the theorems by Narasimhan and Seshadri \parencite{NS1965} and Ramanathan \parencite{Ramanathan1996}. This section relies heavily on the concepts introduced in appendix \ref{appendix_non_connected_lie_groups}. We will not give a detailed proof of the non-connected version of the correspondence, but we will see how it can be reduced to a statement concerning a twisted principal bundle over a covering of the original manifold. We will not give a proof of the latter. To sum up, what this section intends to do is to formulate a version of the Hitchin-Kobayashi for bundles with non-connected structure groups in terms of an equivalent statement concerning a twisted principal bundle on a covering. We sketch a proof of the fact that these two statements are equivalent. 

For the whole section we fix a complex reductive Lie group \(\hat{G}\), and \(\hat{K}\subset \hat{G}\) a maximal compact subgroup. \(K\) and \(G\) will be their respective identity components. We also set \(\Gamma:=\hat{G}/G  \cong\hat{K}/K\) the group of connected components. Then, \(K\) is a maximal compact subgroup of \(G\). We fix \(\langle \cdot , \cdot \rangle\) a non-degenerate \(\C\)-bilinear pairing on \(\mathfrak{g}\) invariant under the adjoint \(\hat{G}\) action. If \(\hat{G}\) is semisimple, then the Killing form \(\kappa_{\mathfrak{g}}\) does the job. If \(\hat{G}\) is only reductive, an invariant pairing still exists \parencite{Benito2024}. Indeed, reductive Lie algebras can be decomposed as \(\mathfrak{g}=\mathfrak{z}(\mathfrak{g})\oplus [\mathfrak{g}, \mathfrak{g}]\). On the semisimple part, the Killing form gives a \(\C\)-bilinear, \(\hat{G}\)-invariant, non-degenerate pairing. On the center, we may choose \textit{any} non-degenerate bilinear form, and it will be invariant under the adjoint action because the center is abelian. The direct sum of these forms yields the desired pairing. If a \(\hat{G}\)-invariant pairing exists, we can make it invariant under the \(\Gamma\)-action on \(\mathfrak{g}\) (cf. equation \eqref{eq_Gamma_action_G}) through a standard averaging procedure, that is, we define a new \(\C\)-bilinear pairing:
\[
\langle \cdot, \cdot \rangle':=\frac{1}{|\Gamma|} \sum_{\gamma \in \Gamma} \langle \cdot \gamma, \cdot \gamma \rangle.
\]
Observe that \(\Gamma\) is finite, hence so is the sum. Recall from the appendix that the action of \(\Gamma\) on \(G\) and \(K\) is given by a lift of the characteristic homomorphism \(\bar{\theta}\colon \Gamma \to \mathrm{Out}(K)\), and its extension to \(G\). In particular, \(K\) is invariant under the \(\Gamma\) action. Since the differential of any Lie group homomorphism is a Lie algebra homomorphism, the induced action on \(\mathfrak{g}\) leaves the Lie algebra of \(Z(K)\), \(\mathfrak{z}_K\), invariant. Indeed, given \(\zeta \in \mathfrak{z}_K, \xi \in \mathfrak{k}\), we have
\begin{equation} \label{lie_group_hom_induces_lie_alg_hom}
    \big[\diff_e\theta_{\gamma}(\zeta), \xi \big] = \diff_e \theta_{\gamma^{-1}} \parent{\big[\zeta, \diff_e \theta_\gamma(\xi) \big]} = \diff_e\theta_{\gamma}(0)=0.
\end{equation}

\subsection{The hermitian Yang-Mills equation revisited}
We saw in section \ref{section_HYM_eq_non_connected} how the hermitian Yang-Mills equation should be generalized to principal bundles with non-connected structure groups. Let \(\mathcal{P}_{\hat{G}}\to X \) be a holomorphic principal \(\hat{G}\)-bundle, and consider a reduction \(P_{\hat{K}}\) of the structure group to \(\hat{K}\). We say that \(P_{\hat{K}}\) is an \textit{\(a\)-hermitian Yang-Mills reduction of the structure group} for \(a \in \Gamma(X;P_{\hat{K}} \times_{\hat{K}}\ii\mathfrak{z}_K)\) satisfying \(\diff_A a=0\) if the Chern connection of the reduction satisfies the \(a\)-hermitian Yang-Mills equation, that is
\begin{equation} \label{eq_HYM_non_connected}
    \Lambda_\omega F_A=- \ii a.
\end{equation}
Equation \eqref{eq_HYM_non_connected} should be understood as an equality between sections of the vector bundle \(P_{\hat{K}}\times_{\hat{K}}\mathfrak{k}\) (observe that it contains \(P_{\hat{K}}\times_{\hat{K}}\mathfrak{z}_K\) as a vector subbundle).

% \Definition{}{Let \(\mathcal{P}_{\hat{G}}\to X \) be a holomorphic principal \(\hat{G}\)-bundle, and consider a reduction \(P_{\hat{K}}\) of the structure group to \(\hat{K}\). Let \(A\) be the Chern connection \(A\) associated to the reduction and \(a \in \Gamma(X;P_{\hat{K}} \times_{\hat{K}}\ii\mathfrak{z}_K)\) satisfying \(\diff_A a=0\). We say that \(A\) satisfies the \textit{\(a\)-Hermitian-Yang-Mills equation} if 
% \begin{equation} \label{eq_HYM_non_connected}
%     \Lambda_\omega F_A=- \ii a,
% \end{equation}
% where \(F_A \in \Omega^2_\R(X; P_{\hat{K}} \times _{\hat{K}} \mathfrak{k})\) is the curvature of the Chern connection.
% }
% Equation \eqref{eq_HYM_non_connected} should be understood as an equality between sections of the vector bundle \(P_{\hat{K}}\times_{\hat{K}}\mathfrak{k}\) (observe that it contains \(P_{\hat{K}}\times_{\hat{K}}\mathfrak{z}_K\) as a vector subbundle).
Sometimes, equation \eqref{eq_HYM_non_connected} can be written in a similar manner to equation \eqref{eq_Hermite_Yang_Mills}, by which we mean that it may be written with a \textit{constant section} on the right-hand side of the equation even if \(\hat{G}\) is \textit{not connected}. This is related to the concept of \textit{monodromy}, which we present briefly. Consider \(p\colon Y \to X\) a Galois \(\Gamma\)-covering between topological spaces, and suppose that \(X\) is path-connected. Fix a basepoint \(x \in X\). The fundamental group \(\pi_1(X,x)\) acts on \(p^{-1}(x)\) (which is finite) from the right as follows: for every equivalence class of loops \([\beta]\in \pi_1(X,x)\), and \(y \in p^{-1}(x)\), we define \(y\cdot [\beta]:=\tilde{\beta}_y(1)\), where \(\tilde{\beta}\) is the unique lift of any representative \(\beta\) of \([\beta]\) whose value at zero is \(y\). Since \(\tilde{\beta}\) is a lift, \(\beta(1)\in p^{-1}(x)\), and everything is well-defined (see \parencite[Th. 11.22]{lee_introTM}). This is what we call the \textit{monodromy action}. The \textit{monodromy} of the covering is the morphism of groups \(\pi_1(X,x)\to \mathrm{Sym}(Y_0)\) where \(Y_0:=p^{-1}(x)\) and \(\mathrm{Sym}(Y_0)\) is the group of permutations of its elements. The image \(M\) of this morphism is what we call the \textit{monodromy group}. Different choices of the basepoint give rise to isomorphic monodromy groups. Finally, whenever we have a principal \(\hat{G}\)-bundle \(P_{\hat{G}}\to X\), we define its monodromy as the monodromy of the covering \(Y:=P_{\hat{G}}/G \to X\). The covering space \(Y\) is connected if and only if the monodromy is surjective.

\begin{obs} \label{obs_HK_non_connected_alpha_a}
   If \(\hat{G}\) is not connected but \(\mathcal{P}_{\hat{G}}\) has  surjective monodromy, then the section \(a\in \Gamma(X;P_{\hat{K}}\times _{\hat{K}} \ii \mathfrak{z}_K)\) can be identified with a single element \(\alpha\in \ii \mathfrak{z}_K^\Gamma\). Then, solving \eqref{eq_HYM_non_connected} is equivalent to solving \(\Lambda_\omega F_A=-\ii \alpha\). To see why, first observe that, since \(\hat{G}\) acts on \(P_{\hat{G}}\), proposition \ref{prop_gamma_Action_Y} implies that there is a natural \(\Gamma\)-action on \(Y:=P_{\hat{G}}/G\). In fact, \(Y \to X\) has the structure of a principal \(\Gamma\)-bundle. We thus have a natural vector bundle isomorphism given by 
    \[
        P_{\hat{K}} \times_{\hat{K}} \mathfrak{z}_K \longrightarrow Y \times_\Gamma \mathfrak{z}_K, \ [p, \zeta] \mapsto [\hat{\pi}(p), \zeta] \
    \]
    where \(\hat{\pi}\colon P_{\hat{G}}\to Y\) is the natural projection (which we restrict to \(P_{\hat{K}}\)). Moreover, the fact that \(\diff_A \alpha=0\) means that the section \(a\) is parallel along every curve in \(X\). What this means, in particular, is that the image of \(a\) at any point is determined uniquely by its image at one point by parallel transport (see \parencite[Prop. 7.6]{lee_introCM}). Now, fix any \(x_0\in X\), and let \(Y_0\) be as in the comments above. \(a\) descends to a \(\Gamma\)-equivariant map \(\delta_a \colon Y_0 \to \mathfrak{z}_K\) given by \(\delta_a(y_0)=\xi_0\) if and only if \(a(x_0)=[y_0, \xi_0]\) (in fact, \(\delta_a\) can be defined in all \(Y\) using the same argument). Since \(Y_0\) is a \(\Gamma\)-orbit, and \([y_0\cdot \gamma, \xi_0\cdot \gamma]=[y_0, \xi_0]\) we see that the map is indeed \(\Gamma\)-equivariant. If we show that \(\delta_a\) is constant, its image must be \(\Gamma\)-invariant, that is, an element of \(\mathfrak{z}_K^\Gamma\). Hence, we will be able to identify \(\delta_a\) with some \(\alpha\in \mathfrak{z}_K^\Gamma\) and conversely, given \(\alpha\) we will be able to recover \(a\) using parallel transport. The map \(\delta_a\) descends to a \(\Gamma\)-equivariant map \(M\setminus Y_0\to \mathfrak{z}_K\). This follows from the fact that \(\diff_A a=0\) forces \(\delta_a\) to be constant on each connected component of \(Y\), and the quotient \(M \setminus Y_0\) parametrizes connected components of \(Y_0\). Since the monodromy action is defined by liftings of paths of the fundamental group, \(y\) and \(y\cdot [\beta]\) are on the same connected component. Hence, \(\delta_a\) takes the same value on both of them. If the monodromy is surjective, then \(M=\mathrm{Sym}(Y_0)\), so \(M\setminus Y_0\) is a point. Then the action of \(\Gamma\) is trivial, and since the map is \(\Gamma\)-equivariant, its image must be an element \(\alpha \in \ii \mathfrak{z}_K\) that is \(\Gamma\)-invariant. That is \(\alpha \in \ii \mathfrak{z}_K^\Gamma\).
    
    A reduction of the structure group to \(\hat{K}\) satisfying \(\Lambda_\omega F_A=-\ii \alpha\) is called an \(\alpha\)\textit{-hermitian Yang-Mills reduction of the structure group}. Observe that \(\alpha\) defines a constant section. Indeed, for every \(\hat{k}\in \hat{K}\), there is \(\gamma \in \Gamma\) such that \(\Ad{\hat{k}}(\alpha)=\diff_e \theta_{\gamma}(\alpha)=\alpha\).
\end{obs}

\subsection{Stability condition}
The first step towards the generalization of the correspondence is to give a notion of stability. In fact, we will need two notions: one for principal bundles with non-connected structure group, and another for twisted equivariant bundles. 

Observe that the statement of lemma \ref{lemma_character} holds in the non-connected case. Moreover, if the pairing \(\langle \cdot, \cdot \rangle\) is \(\Gamma\)-invariant, and \(\zeta \in \ii \mathfrak{k}^\Gamma\), then \(\chi_\zeta\) is \(\Gamma\)-invariant, i.e. \(\chi_\zeta(\xi \cdot \gamma)=\chi_\zeta(\xi)\) for all \(\xi \in \mathfrak{g}\). Note as well that the definition of the degree of a holomorphic principal bundle with respect to a reduction and an element of \(\ii \mathfrak{k}\) applies verbatim. 

Before giving the definition of stability, observe that if we choose \(\zeta \in \ii \mathfrak{k}^\Gamma\), then the parabolic subgroup \(P_\zeta\) defined by proposition \ref{prop_charact_parabolic_subgroups} is a \(\Gamma\)-invariant subgroup of \(G\) (see \parencite[Prop. 2.12]{HKnonconnected}). From proposition \ref{prop_GquotH_Gamma_action} and the comments below, we deduce that \(G/P_\zeta\) is a \((\theta, c)\)-twisted \((G, \Gamma)\)-equivariant manifold. Finally, by proposition \ref{prop_gamma_equivariant_bundles}, we deduce that, given a \((\theta,c)\)-twisted \(\Gamma\)-equivariant bundle, \(P_G\to Y\), the associated bundle \(P_G(G/P_\zeta)\to Y\) is a \(\Gamma\)-equivariant principal bundle. Hence, we may consider \(\Gamma\)-equivariant sections or, equivalently, \(\Gamma\)-equivariant reductions of the structure group to \(P_\zeta\). Similarly, consider \(L_\zeta \subset P_\zeta\). It is \(\Gamma\)-invariant, and given a \(\Gamma\)-invariant reduction of the structure group of \(P\) to \(P_\zeta\), call it \(P_\sigma\), we know from proposition \ref{prop_Gamma_inv_red_struc_group} that \(P_\sigma\) is endowed with the structure of a \((\theta,c)\)-twisted \(\Gamma\)-equivariant principal \(P_\zeta\)-bundle. Since \(L_\zeta\subset P_\zeta\) is also \(\Gamma\)-invariant, and \(\hat{L}_\zeta \subset \hat{P}_\zeta\) meets all the connected components of the latter, we have an isomorphism \(P_\zeta/L_\zeta\cong \hat{P}_\zeta/\hat{L}_\zeta\). Thus, the same arguments apply, \(P_\sigma(P_\zeta/L_\zeta)\) is a \(\Gamma\)-equivariant principal bundle, and it makes sense to consider \(\Gamma\)-equivariant reductions of the structure group to \(L_\zeta\). With this in mind, we are now ready to define stability for twisted principal bundles. 

\Definition{}{Let \(\mathcal{P}\) be a \((\theta,c)\)-twisted \(\Gamma\)-equivariant principal bundle over a Riemann surface \(Y\), and let \(\alpha \in \ii \mathfrak{z}_K^\Gamma\). We say that \(\mathcal{P}\) is
\begin{enumerate}
    \item \textit{\(\Gamma\)-invariantly \(\alpha\)-stable (resp. semistable)}\index{Principal bundle!\(\Gamma\)-invariantly \(\alpha\)-semistable} if for any \(\zeta \in \ii \mathfrak{k}^\Gamma\) and any \(\Gamma\)-invariant holomorphic section \(\sigma \in \mathcal{O}(Y;\mathcal{P}(G/P_\zeta))\), we have 
    \[
    \deg \mathcal{P}(\sigma, \zeta)- \langle \alpha, \zeta\rangle >0 \hspace{1 cm} \parent{\text{resp. } \  \deg \mathcal{P}(\sigma, \zeta)-\langle \alpha, \zeta\rangle \geq 0}.
    \]
    \item \textit{\(\alpha\)-invariantly \(\Gamma\)-polystable}\index{Principal bundle!\(\Gamma\)-invariantly \(\alpha\)-polystable} if it is \(\Gamma\)-invariantly \(\alpha\)-semistable and, for any \(\zeta \in \ii \mathfrak{k}^\Gamma\) and any \(\Gamma\)-invariant holomorphic section \(\sigma \in \mathcal{O}(Y;\mathcal{P}(G/P_\zeta))\) such that \(\deg \mathcal{P}(\sigma, \zeta) - \langle \alpha, \zeta \rangle = 0 \), there is a \(\Gamma\)-invariant holomorphic reduction of the structure group \(\sigma'\in \mathcal{O}(X;\mathcal{P}_\sigma(P_\zeta/L_\zeta))\).
\end{enumerate}
}

Now we need a notion of stability for principal bundles with non-connected structure group. The definition we are about to give is quite similar to the definition for twisted equivariant principal bundles, and as we shall see, these definitions are closely related. 

\Definition{}{Let \(X\) be a compact Riemann surface, and let \(\mathcal{P}(X,\hat{G})\) be a holomorphic principal \(\hat{G}\)-bundle with surjective monodromy, and let \(\alpha \in \ii \mathfrak{z}_K^\Gamma\). We say that the principal bundle is
\begin{enumerate}
    \item \textit{\(\alpha\)-stable (resp. semistable)}\index{Principal bundle!\(\alpha\)-semistable} if for any \(\zeta \in \ii \mathfrak{k}^\Gamma\) and any holomorphic section \(\sigma \in \mathcal{O}(X; \mathcal{P}(\hat{G}/\hat{P}_\zeta))\), we have
    \[
    \deg \mathcal{P}(\sigma, \zeta) - \langle \alpha, \zeta\rangle >0 \hspace{1 cm} \parent{\text{resp.} \ \ \deg \mathcal{P}(\sigma, \zeta)-\langle \alpha, \zeta\rangle \geq 0}
    \]
    \item \textit{\(\alpha\)-polystable}\index{Principal bundle!\(\Gamma\)-invariantly \(\alpha\)-stable} if it is \(\alpha\)-semistable and for any \(\zeta \in \ii \mathfrak{k}^\Gamma\) and \(\sigma \in \mathcal{O}(X; \mathcal{P}(\hat{G}/\hat{P}_\zeta))\) such that \(\deg \mathcal{P}(\sigma, \zeta)-\langle \alpha, \zeta \rangle=0\), there is a holomorphic reduction of the structure group \(\sigma'\in \mathcal{O}(X; \mathcal{P}_\sigma(\hat{P}_\zeta/\hat{L}_\zeta))\).
\end{enumerate}
}

When \(G=\hat{G}\) is connected, we recover the usual stability criteria (see definition \ref{def_stability_connected}).

\subsection{The theorem}
The aim of this section is to state the version of the correspondence for principal bundles with non-connected structure group and sketch how it can be reduced to a \(\Gamma\)-equivariant correspondence. We fix a complex reductive non-connected Lie group \(\hat{G}\), its identity component \(G\), \(\hat{K}\) a maximal compact subgroup, and \(K:=\hat{K}\cap G\) its identity component. We thus have the following diagram of extensions
\begin{center}
     \begin{tikzcd}[row sep=normal, column sep=normal, ampersand replacement=\&]
     1 \arrow[rr] \& \& G \arrow[rr] \& \& \hat{G} \arrow[rr] \& \& \Gamma \arrow[rr] \& \&  1 \\
     1 \arrow[rr] \& \& K \arrow[hookrightarrow, u] \arrow[rr] \& \& \hat{K}\arrow[hookrightarrow, u] \arrow[rr] \& \& \Gamma \ar[u,-,double equal sign distance,double]  \arrow[rr]\& \&  1
     \end{tikzcd}
\end{center}
We also fix a Riemann surface \(X\), and a lift of the characteristic homomorphism \(\theta\colon \Gamma\to \mathrm{Aut}(K)\) (which we extend to a lift to \(\mathrm{Aut}(G)\)), and a two-cocycle \(c\colon \Gamma \times \Gamma \to Z(K)\) as in appendix \ref{appendix_lie}. Under these hypotheses, the version of the Hitchin-Kobayashi correspondence for principal bundles with non-connected structure group reads as follows.

\Theorem{\textbf{Hitchin-Kobayashi, non-connected version}}{Let \(\alpha\in \ii \mathfrak{z}_K^\Gamma\). A holomorphic principal \(\hat{G}\)-bundle \(\mathcal{P}_{\hat{G}}\) admits a reduction of the structure group to \(\hat{K}\), say \(P_{\hat{K}}\), satisfying equation \eqref{eq_HYM_non_connected} for a section \(a \in \Gamma(X;P_{\hat{K}} \times_K \mathfrak{z}_K )\) satisfying \(\diff_A a=0\), where \(A\) is the Chern connection in \(\hat{K}\) if and only if it is \(\alpha\)-polystable for \(\alpha \in \ii \mathfrak{z}_K\) defined as in remark \ref{obs_HK_non_connected_alpha_a}. \label{th_HK_non_connected}}

Our strategy for the proof is to translate the problem into one concerning principal bundles with connected structure groups. This is done by means of the correspondence of categories given by theorem \ref{th_equivalence_categories}. More explicitly, given a holomorphic principal \(\hat{G}\)-bundle over \(X\), \(\mathcal{P}_{\hat{G}}\), we define \(Y:=\mathcal{P}_{\hat{G}}/G\), and consider the quotient map \(\mathcal{P}_{\hat{G}}\to Y\), which endows \(\mathcal{P}_{\hat{G}}\) with the structure of a \((\theta, c)\)-twisted \(\Gamma\)-equivariant principal \(G\)-bundle over \(Y\). We denote it \(\mathcal{P}_G\), even though, as a set, and as a manifold, it coincides with \(\mathcal{P}_{\hat{G}}\). We will suppose that the monodromy of the covering is surjective, so that \(Y\) is connected and remark \ref{obs_HK_non_connected_alpha_a} applies. If the monodromy is not surjective, one may consider intermediate reductions of the structure group to make it connected, but we take the surjectivity as an assumption and not delve deeper into the issue.

To prove theorem \ref{th_HK_non_connected}, the idea is to show that it is equivalent to the following theorem.

\Theorem{\textbf{Hitchin-Kobayashi for twisted bundles}}{Let \(P_G\to X\) be a \((\theta,c)\)-twisted \(\Gamma\)-equivariant principal \(G\)-bundle, and let \(\alpha \in \ii \mathfrak{z}_K^\Gamma\). Then, \(P\) is \(\Gamma\)-invariantly \(\alpha\)-polystable if and only if there exists a \(\Gamma\)-equivariant reduction of the structure group \(\sigma \in \Gamma(X; P_G(G/K))\) satisfying \(\Lambda_\omega F_A=-\ii \alpha\), where \(A\) is the Chern connection on \(P_K\). \label{th_HK_twisted_bundles}}
A \(\Gamma\)-equivariant reduction of the structure group satisfying \(\Lambda_\omega F_A=-\ii \alpha\) is called a \(\Gamma\)\textit{-equivariant} \(\alpha\)\textit{-hermitian Yang-Mills reduction of the structure group} to \(K\). Moreover, observe that a \(\Gamma\)-invariant reduction of the structure group is, in particular, a reduction of the structure group, hence it makes sense to consider its Chern connection.

We are not claiming that theorem \ref{th_HK_twisted_bundles} is straightforward, or a direct consequence of the Hitchin-Kobayashi correspondence for principal bundles with connected structure group because we have to take into account the twisted structure somehow. However, since we are now working with a principal bundle with connected structure group over a connected Riemann surface, it seems reasonable to think that a proof in terms of a functional \(\Psi\), as we did in section \ref{section_sketch_proof} will work here as well.

It remains to sketch how theorem \ref{th_HK_non_connected} can be reduced to theorem \ref{th_HK_twisted_bundles}. The first thing one should check is that the condition of \(\alpha\)-polystability for the bundle \(\mathcal{P}_{\hat{G}}\to X\) is equivalent to the condition of \(\alpha\)-polystability of the twisted equivariant bundle \(P_G \to Y\). More precisely:

\Proposition{}{Let \(\mathcal{P}_{\hat{G}}\) be a holomorphic principal \(\hat{G}\)-bundle, and let \(\mathcal{P}_G\) be the corresponding \((\theta,c)\)-twisted \(\Gamma\)-equivariant principal bundle over \(Y:=\mathcal{P}_{\hat{G}}\ /G\) given by theorem \ref{th_equivalence_categories}. Let \(\alpha \in \ii \mathfrak{z}_K^\Gamma\), then \(\mathcal{P}_{\hat{G}}\) is \(\alpha\)-polystable if and only if \(\mathcal{P}_G\) is \(\Gamma\)-invariantly \(\alpha\)-polystable. \label{prop_equiv_twisted_non_con}}
\begin{proof}[\textit{Sketch of the proof}]
    Fix \(\zeta \in \ii \mathfrak{k}\). One may show that \(G/P_\zeta \cong \hat{G}/\hat{P}_\zeta\) and that there is a bijection between holomorphic sections of the bundle \(\mathcal{P}_{\hat{G}}(\hat{G}/\hat{P}_\zeta)\to X\) (that is, holomorphic reductions of the structure group to \(\hat{P}_\zeta\)), and \(\Gamma\)-equivariant sections of the bundle \(\mathcal{P}_G(G/P_\zeta) \to Y\) (that is, \(\Gamma\)-equivariant reductions of the structure group to \(P_\zeta\)). Call \(\sigma\) the section, under some abuse of notation, we will use the same notation for a section seen from both sides of the previous bijection. A similar discussion applies to the Levi subgroup, which yields a bijection between holomorphic sections of \((\mathcal{P}_{\hat{G}})_\sigma(\hat{P}_\zeta/\hat{L}_\zeta)\to X\) and \(\Gamma\)-equivariant sections of \((\mathcal{P}_G)_\sigma(P_\zeta/L_\zeta)\to Y\). This reasoning is used to show that \(\deg \mathcal{P}_G(\sigma, \zeta)=\deg \mathcal{P}_{\hat{G}}(\sigma, \zeta)\), and from this identity the claim follows directly.
\end{proof}

Finally, one shows that \(\mathcal{P}_{\hat{G}}\to X\) admits an \(\alpha\)-hermitian Yang-Mills reduction of the structure group to \(\hat{K}\) if and only if \(P_G\to Y\) admits a \(\Gamma\)-equivariant \(\alpha\)-hermitian Yang-Mills reduction of the structure group to \(K\).

% We then show that the statement of the correspondence for non-connected structure group is equivalent to a Hitchin-Kobayashi correspondence for twisted equivariant bundles over a covering, but we do not give a proof of this result. As a final comment before going into the details, whenever we have a holomorphic principal bundle over a compact Kähler manifold \(\mathcal{P}_{\hat{G}}\to X\), we will use theorem \ref{th_equivalence_categories} to define \(Y:=\mathcal{P}_{\hat{G}}/G\), and obtain a twisted equivariant covering \(\mathcal{P}_G \to Y\). We will assume that the monodromy of \(\mathcal{P}_{\hat{G}}\to X\) is surjective. This implies that the covering \(Y\) is connected. If the monodromy is not connected, one may consider intermediate reductions of the structure group to make it connected, but we take the surjectivity as an assumption a not dwelve deeper into the issue.

% \Proposition{}{Let \(\mathcal{P}_{\hat{G}}\) be a holomorphic principal \(\hat{G}\)-bundle, and let \(\mathcal{P}_G\) the corresponding \((\theta,c)\)-twisted \(\Gamma\)-equivariant principal bundle over \(Y:=\mathcal{P}_{\hat{G}}\ /G\) given by theorem \ref{th_equivalence_categories}. Let \(\alpha \in \ii \mathfrak{z}_K^\Gamma\), then \(\mathcal{P}_{\hat{G}}\) is \(\alpha\)-polystable if and only if \(\mathcal{P}_G\) is \(\Gamma\)-invariantly \(\alpha\)-polystable. \label{prop_equiv_twisted_non_con}}
% Before given the proof, we will need some lemmas.
% \Lemma{}{Let \(\hat{G}\) be an \(\mathcal{S}\)-Lie group, \(G\) its identity component, and \(\hat{H} \subset \hat{G}\) a closed embedded \(\mathcal{S}\)-Lie subgroup that meets every connected component of \(\hat{G}\). Let \(H:=\hat{H}\cap G\). Then, \(\hat{G}/\hat{H}\) is \(\mathcal{S}\)-isomorphic to \(G/H\). \label{lemma_HK_non_connected_1}}
% \begin{proof}
%     Define \(G/H \to \hat{G}/\hat{H}, \ gH \mapsto g \hat{H}\). This map is well-defined: if \(gH=g'H\), then \(g^{-1}g'\in H\subset \hat{H}\), hence \(g\hat{H}=g'\hat{H}\). It is also injective: if \(g\hat{H}=g'\hat{H}\), then \(g^{-1}g'\in \hat{H}\cap G=H\), and surjective. To see surjectivity, observe that all the connected components of \(\hat{G}\) are given by the cosets \(\hat{g} G\) for \(g \in \hat{G}\). Let \(\hat{g}\hat{H}\in \hat{G}/\hat{H}\). We want to see that there is \(g \in G\) such that \(\hat{g}\hat{H}=g\hat{H}\). The connected component containing \(\hat{g}\) is \(\hat{g}G\). Since \(\hat{H}\) meets every connected component of \(\hat{G}\), \(\hat{g}G \cap \hat{H} \neq \emptyset\), hence there is \(h \in H\) such that \(\hat{g}g=h\). Since \(G\triangleleft \hat{G}\), \(\hat{g}G=G\hat{g}\), so there is \(g_2\in G\) such that \(\hat{g}g=g_2\hat{g}=h\). Choosing \(g:=g_2^{-1}\), we have \(g\hat{H}=g_2^{-1}h \hat{H}=\hat{g}H\). Finally, the map is an \(\mathcal{S}\)-isomorphism by the characteristic property of surjective submersions and the equivariant rank theorem \parencite[Cor. 2.6]{lee_introCM} and \parencite[Th. 7.25]{lee_introSM}.
% \end{proof}

% \Lemma{}{Let \(F\) be a \((\theta,c)\)-twisted \((G, \Gamma)\)-manifold and \(\pi\colon P \to X\) be a \((\theta, c)\)-twisted \(\Gamma\)-equivariant principal \(G\)-bundle. Let \(\mathscr{C}^\infty(P, F)^{G, \Gamma}\) the space of \((\theta,c)\)-twisted \((G, \Gamma)\)-equivariant smooth maps \eqref{eq_def_twisted_equiv_map}, and let \(\Gamma_{\mathrm{inv}}(X, P(F))\) the space of \(\Gamma\)-invariant sections on the twisted equivariant bundle \(P \to X\). Then, there is a bijection \(\mathscr{C}^\infty(P, F)^{G, \Gamma} \leftrightarrow \Gamma_{\mathrm{inv}}(X, P(F))\). The result holds form holomorphic sections and maps as well. \label{lemma_HK_non_connected_2}}
% \begin{proof}
%     See \parencite[Prop. 4.8]{barajas2023non}. Schematically, given a section \(s\in\Gamma_{\text{inv}}(X;P(F))\) we define \(\tilde{s}\in \mathscr{C}^\infty(P, F)^{G, \Gamma}\) by \(\tilde{s}(p)=f\) if \(s([p])=[p,f]\) where we have identified \(X\cong P/G\). Conversely, given a map \(\tilde{s}\in \mathscr{C}^\infty(X, P(F))\) we define \(s([p])=[p, \tilde{s}(p)]\).
% \end{proof}

% \begin{proof}[\textit{Proof of proposition} \ref{prop_equiv_twisted_non_con}]
%     First, observe that on both definitions of stability, we ask for \(\alpha \in \ii \mathfrak{z}_K^\Gamma\), and we choose elements \(\zeta \in \ii \mathfrak{k}^\Gamma\). Moreover, given \(\zeta \in \ii \mathfrak{k}^\Gamma\), the group \(\hat{P}_\zeta\) meets all the connected components of \(\hat{G}\) \parencite[Prop. 2.12]{HKnonconnected}. Since \(P_\zeta=\hat{P}_\zeta \cap G\), lemma \ref{lemma_HK_non_connected_1} implies that \(G/P_\zeta=\hat{G}/\hat{P}_\zeta\). This implies that any holomorphic section of the associated bundle \(\mathcal{P}_{\hat{G}}(\hat{G}/\hat{P}_\zeta)\to X\) induces a \(\Gamma\)-invariant section of the \(\Gamma\)-equivariant bundle \(\mathcal{P}_{G}(G/P_\zeta)\to Y\) and vice versa. To see this, first observe that \(\mathcal{P}_{\hat{G}}=\mathcal{P}_G\) as manifolds (cf. remark \ref{obs_equivalence_categories}). Observe that \(G/P_\zeta\cong \hat{G}/\hat{P}_\zeta\) is a \((\theta,c)\)-twisted \((G, \Gamma)\)-equivariant manifold (cf. proposition \ref{prop_GquotH_Gamma_action} and the comments below). Moreover, a map between twisted equivariant manifolds is \((\theta,c)\)-twisted \((G, \Gamma)\)-equivariant if and only if it is \(\hat{G}\)-equivariant. Sections of \(\mathcal{P}_{\hat{G}}(\hat{G}/\hat{P}_\zeta)\to X\) are in one-to-one correspondence with \(\hat{G}\)-equivariant maps \(\mathcal{P}_{\hat{G}}\to G/P_\zeta\), and these are in one-to-one correspondence with \((\theta,c)\)-twisted \((G, \Gamma)\)-equivariant maps \(\mathcal{P}_{\hat{G}}\to \hat{G}/\hat{P}_\zeta\), observe that these coincide with \((\theta,c)\)-twisted \((G, \Gamma)\)-equivariant maps \(\mathcal{P}_{{G}}\to G/P_\zeta\) because the sets are the same. Moreover, by proposition \ref{prop_gamma_equivariant_bundles} we know that \(\mathcal{P}_{\hat{G}}(\hat{G}/\hat{P}_\zeta)\to X\) is a \(\Gamma\)-equivariant fiber bundle. Applying lemma \ref{lemma_HK_non_connected_2}, to each \((\theta,c)\)-twisted \((G, \Gamma)\)-equivariant map \(\mathcal{P}_{{G}}\to G/P_\zeta\) corresponds a \(\Gamma\)-invariant section of the \(\Gamma\)-equivariant bundle \(\mathcal{P}_G \to Y\). Under some abuse of notation, we shall use the same notation for both sections. Observe that by proposition \ref{prop_Gamma_inv_red_struc_group} this means that there is a one-to-one correspondence between holomorphic reductions of the structure group of \(\mathcal{P}_{\hat{G}}\to X\) to a parabolic subgroup \(\hat{P}_\zeta\), and \(\Gamma\)-invariant reductions of the twisted equivariant bundle \(\mathcal{P}_G \to Y\) to the connected component \(P_\zeta\) of \(\hat{P}_\zeta\). 

%     The same discussion applies to Levi subgroups: \(\hat{L}_\zeta\) meets every connected component of \(\hat{P}_\zeta\), \(P_\zeta\) is the identity component of \(\hat{P}_\zeta\), and \(L_\zeta =\hat{L}_\zeta \cap P_\zeta\). Hence, lemma \ref{lemma_HK_non_connected_1} holds, and \(\hat{P}_\zeta /\hat{L}_\zeta\cong P_\zeta/L_\zeta\). By the same reasoning, a section of the bundle \((\mathcal{P}_{\hat{G}})_\sigma(\hat{P}_\zeta / \hat{L}_\zeta)\to X\) is equivalent to a \(\Gamma\)-invariant section of the \(\Gamma\)-equivariant bundle \((\mathcal{P}_G)_\sigma(P_\zeta/L_\zeta)\to Y\). Observe as well that the \(\Gamma\)-invariance of \(\zeta\) ensures that the integral over \(Y\) descends to an integral over \(X\). In particular, we have \(\deg \mathcal{P}_{\hat{G}}(\sigma, \zeta)=\deg \mathcal{P}_G(\sigma, \zeta)\).
    
%     {\color{red}Aquí en el paper solo afirmáis que la integral desciende a una integral en \(X\) porque siempre se puede encontrar una conexión \(A\) que sea \(\Gamma\)-invariante. Esto no lo entiendo. Primero, aunque \(Y/\Gamma=X\), no tengo muy claro como se relacionan las integrales entre variedades y sus cocientes. He encontrado algo llamado ``integración a lo largo de fibras'' que es una especie de Fubini para variedades cociente, aunque no sé si es lo que estamos aplicando aquí. Según ese concepto de integración a lo largo de fibras, si tenemos un recubrimiento, es como un \(\Gamma\)-fibrado, entonces para las formas \(\alpha\) en \(Y\) que son invariantes, existen formas \(\beta\) en \(X\) de manera que 
%     \[
%     \int_Y \alpha = |\Gamma|\int_X \beta
%     \]    
%     Consideramos el fibrado twisted \(P_G \to Y\). Suponiendo que \(K_\zeta\) sea \(\Gamma\)-invariante tiene sentido hablar de una reducción \((P_G)'_\sigma\) que sea \(\Gamma\)-invariante ?. Entonces si la 2-forma \(\alpha=\chi_\zeta(F_A)\) se correspondiera con una 2-forma \(\chi_\zeta(F_{\hat{A}})\) para alguna conexión \(\hat{A}\) en el fibrado \(P_{\hat{G}}\to X\), tendríamos 
%     \[
%     \int_Y \chi_\zeta(F_{A})=|\Gamma|\cdot \int_X \chi_\zeta(F_{\hat{A}}).
%     \]
%     Esto sigue sin ser lo que queremos porque aparece un factor de proporcionalidad \(|\Gamma|\), pero si de alguna forma nos pudiéramos quitar el factor de proporcionalidad sí que tendríamos igualdad de los grados, que entiendo que es lo que queremos. En cualquier caso, no sé si con la frase de que la integral desciende a una en \(X\) os referís a esto o a otra cosa.
%     } 

%     Having said that, suppose that \(\mathcal{P}_{\hat{G}}\) is \(\alpha\)-polystable for \(\alpha \in \ii \mathfrak{z}_K^\Gamma\), and let \(\zeta \in \ii \mathfrak{k}^\Gamma\). Let \(\sigma \in \mathcal{O}(Y; \mathcal{P}_{{G}}(G/P_\zeta))\) be \(\Gamma\)-invariant. We have seen that this induces a section (which we cal the same) \(\sigma \in \mathcal{O}(X; \mathcal{P}_{\hat{G}}(\hat{G}/\hat{P}_\zeta))\). Now, since \(\mathcal{P}_{\hat{G}}\) is \(\alpha\)-polystable, it is \(\alpha\)-semistable, hence 
%     \begin{equation} \label{eq_P_G_1}
%     \deg \mathcal{P}_{\hat{G}}(\sigma, \zeta)-\langle \alpha, \zeta\rangle=\deg \mathcal{P}_{{G}}(\sigma, \zeta)-\langle \alpha, \zeta \rangle \geq 0,
%     \end{equation}
%     so \(\mathcal{P}_{G}\) is \(\Gamma\)-invariantly \(\alpha\)-semistable. Now, take \(\zeta \in \ii \mathfrak{k}^\Gamma\) and a \(\Gamma\)-invariant \(\sigma \in \mathcal{O}(Y; \mathcal{P}_G(G/P_\zeta))\) such that the equality holds in \eqref{eq_P_G_1}. Again, this induces a section \(\sigma\in \mathcal{O}(X; \mathcal{P}_{\hat{G}}(\hat{G}/\hat{P}_\zeta))\) such that the equality holds. Since \(\mathcal{P}_{\hat{G}}\) is \(\alpha\)-polystable, there is a holomorphic section \(\sigma'\in \mathcal{O}(X; \mathcal{P}_{\hat{G}}(\hat{P}_\zeta/\hat{L}_\zeta))\) which induces a \(\Gamma\)-invariant section (again, we keep the name) \(\sigma'\in \mathcal{O}(Y; \mathcal{P}_G(P_\zeta/L_\zeta))\). The converse follows from the same reasoning in the opposite direction.
% \end{proof}

% \Theorem{\textbf{Hitchin-Kobayashi for twisted bundles}}{Let \(P\to X\) a \((\theta,c)\)-twisted \(\Gamma\)-equivariant principal \(G\)-bundle, and let \(\alpha \in \ii \mathfrak{z}_K^\Gamma\). Then, \(P\) is \(\Gamma\)-invariantly \(\alpha\)-polystable if and only if there exists a \(\Gamma\)-equivariant reduction of the structure group \(\sigma \in \Gamma(X; P(G/K))\) satisfying \(\Lambda_\omega F_A=-\ii \alpha\), where \(A\) is the Chern connection on \(P_K\). \label{th_HK_twisted_bundles}}
% A \(\Gamma\)-equivariant reduction of the structure group satifying \(\Lambda_\omega F_A=-\ii \alpha\) is called a \(\Gamma\)\textit{-equivariant} \(\alpha\)\textit{-hermitian Yang-Mills reduction of the structure group} to \(K\). Moreover, observe that a \(\Gamma\)-invariant reduction of the structure group is, in particular, a reduction of the structure group, hence it makes sense to consider its Chern-connection. The following results are intended to show how this result is equivalent to theorem \ref{th_HK_non_connected}.


% \Proposition{}{Let \(\mathcal{P}_{\hat{G}}\) be a principal \(\hat{G}\)-bundle over a compact Riemann surface \(X\), and consider the associated \((\theta,c)\)-twisted \(\Gamma\)-equivariant principal bundle \(\mathcal{P}_G \to \mathcal{P}_{\hat{G}}/G=:Y\). Then, \(\mathcal{P}_{\hat{G}}\) admits an \(\alpha\)-hermitian Yang-Mills reduction of the structure group to \(\hat{K}\) if and only if \(\mathcal{P}_G\) admits a \(\Gamma\)-equivariant \(\alpha\)-hermitian Yang-Mills reduction of structure group to \(K\). \label{prop_proof_HK_nc}}
% \begin{proof}
%     Since \(\hat{K}\) meets all the connected components of \(\hat{G}\), lemma \ref{lemma_HK_non_connected_1} implies that \(G/K\cong \hat{G}/\hat{K}\), Moreover, \(P_G(G/K)/\Gamma \cong P_{\hat{G}}(\hat{G}/\hat{K})\) implies that there is a bijection between the set of sections of the bundle \(P_{\hat{G}}(G/K)\to X\) and the set of \(\Gamma\)-equivariant sections \(P_G(G/K)\to Y\). Hence, suppose that \(P_{\hat{G}} \to X\) admits an \(\alpha\)-hermitian Yang-Mills reduction of the structure group to \(\hat{K}\). This reduction is given by a section \(\sigma\in \Gamma(X; P_{\hat{G}}(\hat{G}/\hat{K}))\) whose Chern connection satisfies \(\Lambda_\omega F_A=-\ii \alpha\). Through the bijection above, the reduction corresponds to a \(\Gamma\)-equivariant section of the \(\Gamma\)-equivariant bundle \(P_G(G/K)\to Y\). Let \(\tilde{A}\) the Chern connection on the reduction \(P_K\), and let \(\tilde{\omega}=p^*\omega\in \Omega^2_\C(Y)\) be the Kähler form on \(Y\), where \(p\colon Y \to X\) is the covering. We want to see that \(\Lambda_{\tilde{\omega}} F_{\tilde{A}}=-\ii \alpha\).
    
%     {\color{red} no sé como seguir la demostración, más que nada porque no acabo de ver como se relaciona la conexión de Chern en \(P_{\hat{K}}\) con la conexión de Chern en \(P_K\). Desde luego, la \(\Lambda_\omega F_A=-\ii \alpha\) no es \textit{exactamente la misma} en ambos casos y ya está, porque según si estamos en un fibrado o en otro, tenemos (a ambos lados de la igualdad) secciones de \(P_K \to Y\) o de \(P_{\hat{K}}\to X\), que obviamente son fibrados distintos.}
% \end{proof}

% With the previous work, and taking for granted theorem \ref{th_HK_twisted_bundles}, we can give a proof of the non-connected version of the correspondene.
% \begin{proof}[\textit{Proof of theorem} \ref{th_HK_non_connected}]
% First, by proposition \ref{prop_equiv_twisted_non_con}, \(\mathcal{P}\) is \(\alpha\)-polystable if and only if \(P_G\) is \(\Gamma\)-invariantly \(\alpha\)-polystable. By theorem \ref{th_HK_twisted_bundles}, this is equivalent to the existence of a \(\Gamma\)-equivariant reduction of the structure group to \(K\) satisfying \(\Lambda_\omega F_A=-\ii \alpha\), i.e. a \(\Gamma\)-equivariant \(\alpha\)-hermitian Yang-Mills reduction of the structrue group. Finally, by proposition \ref{prop_proof_HK_nc}, this is equivalent to the existence of an \(\alpha\)-hermitian Yang-Mills reduction of the structure group to \(\hat{K}\), i.e. a reduction satisfying \(\Lambda_\omega F_A=-\ii \alpha\).
% \end{proof}

\subsection{Moment map picture for the non-connected version}
This version of the correspondence also has a moment map interpretation. Let \(P_{\hat{G}}\) be a smooth principal \(\hat{G}\)-bundle, and consider a reduction of the structure group \(P_{\hat{K}}\) to a principal \(\hat{K}\)-bundle. By the Chern-Singer correspondence, each \(A \in \mathscr{A}\parent{P_{\hat{K}}}\) defines a holomorphic structure on \(P_{\hat{G}}\). Let \(\mathscr{G}\) be the gauge group of \(P_{\hat{K}}\). Upon completing these spaces with Sobolev norms, \(\mathscr{A}\parent{P_{\hat{K}}}\) and \(\mathscr{G}\) are a Banach manifold and a Banach-Lie group respectively. Then, given a section \(a \in \Gamma(X; P_{\hat{K}} \times_K \mathfrak{z}_K)\) satisfying \(\diff_A a=0\), the action of \(\mathscr{G}\) on \(\mathscr{A}\parent{P_{\hat{K}}}\) given by \eqref{eq_def_action_Phi} is Hamiltonian, and the associated moment map is 
\[
\begin{array}{cccc}
    \mu \colon & \mathscr{A}\parent{P_{\hat{K}}} & \longrightarrow & \mathrm{Lie}(\mathscr{G}) \\
    & A & \mapsto & \Lambda_\omega F_A +\ii a
\end{array}
\]
is a moment map (we are identifying \(\mathrm{Lie}(\mathscr{G})\) and its dual as in remark \ref{obs_mu_lie_sin_dual}). Hence, the moduli space of solutions to the hermitian Yang-Mills equations for \(P_{\hat{G}}\) is given by 
\[
\mu^{-1}(0)/\mathscr{G}.
\]
We could make these constructions rigorous as we did for the connected version.

\section{Conclusions and future work} 
The aim of this work has been to present an account of the Hitchin-Kobayashi correspondence for principal bundles with connected structure groups and to set the bases of a version of the correspondence for principal bundles with non-connected structure groups. To that end, we have worked mostly from a differential-geometric viewpoint. Aside from the preliminaries and prerequisites we have covered in the appendices and chapters \ref{chapter_kempf_Ness} and \ref{chapter_principal_bundles}, we have presented a somewhat detailed and self-contained exposition of the principal bundle version of the correspondence. These are relatively standard results today, and they are part of the shared mathematical knowledge of the research community, and we do not claim any originality of the results presented. However, there are not many sources from which one could try to get acquaintance with the topic aside from the original papers and some scattered lecture notes. We hope this work may be useful for future readers trying to learn the material. We shall remark, however, that we have not discussed the algebraic point of view. Recall that the notion of stability arises from Mumford's GIT, but we have presented all our definitions from a differential-geometric point of view. We have also not covered with many details how these results may be phrased rigorously in terms of moduli spaces (aside from our comments in section \ref{section_inf_dim_Kempf_Ness}). 

Our second objective was to lay the foundations for a proof of a correspondence in the non-connected case. We have not given a complete proof of the result because we do not have a proof of \ref{th_HK_twisted_bundles} and, to our knowledge, this result is not proven in the literature. However, we think that reducing the problem to the twisted version is quite an advancement as the twisted version looks much more manageable. Indeed, it seems quite reasonable to think that the techniques outlined in section \ref{section_sketch_proof} may be applied to obtain a proof of the twisted version. 

To sum up, since the publishing of the theorem by Narasimhan and Seshadri in 1965, the discipline has evolved significantly, and several generalizations of the correspondence have appeared in recent years. We hope that the results we present here will be of interest to improve our understanding of the Hitchin-Kobayashi correspondence.

% {\color{blue}
% \textbf{Problemas/Dudas:}
% \begin{enumerate}
%     \item[\textbf{(1)}] En \parencite{garcia2009hitchin} dice: ``Then, by proposition 2.14 \(\mathrm{aut}^{ss}(E, \varphi)=H^0(E(\mathfrak{z}))\), so \((E, \varphi)\) is simple in the sense of Definition 3.8 in \parencite{bradlow2003relative}''. Pero la definición que se da de simple es de un par \((A, \Phi)\) simple, no de que el propio par principal \((E, \varphi)\) sea simple. No he visto la definición de que el propio par principal sea simple en ningún sitio, así que he asumido que quiere decir que tiene una conexión simple.
%     \item[\textbf{(2)}] Esto directamente no se si es verdad. 
%     \item[\textbf{(3)}] Este enunciado no aparece explícitamente en el paper, intenta ser un resumen de lo que llamáis descomposición de Jordan-Holder dando los mínimos detalles posibles. No sé si estará bien. 
%     \item[\textbf{(4)}] En la sección de espacios de Sobolev (sección \ref{section_Sobolev_spaces}) defino esas normas de una manera concreta, pero no estoy seguro de que esté bien. Primero, en la mayoría de sitios donde he visto esto hecho simplemente escriben 
%     \[
%     \langle \psi. \phi\rangle_k=\int_X \sum_{i=1}^r \langle \nabla ^i \psi, \nabla^i \phi \rangle
%     \]
%     y no dicen que es \(\nabla^i\). En el TFM de Diego he visto que define la norma para formas a valores en un fibrado vectorial cualquiera con una conexión a partir del operador \(\diff_\nabla\) como
%     \[
%     ||\alpha||_k:=\sqrt{\sum_{i=0}^r ||\diff_\nabla ^i \alpha ||_{L_2}^2}.
%     \]
%     No sé si la definición que he dado yo es equivalente a la anterior para espacios de formas. Lo que hago yo lo he visto en el libro \parencite{wehrheim2004uhlenbeck}. He mirado el libro de Donalson \parencite{donaldson1990geometry} y yo diría que hace lo mismo que yo, pero es de esos que no parece decir explícitamente a que se refiere con \(\nabla^i\).
%     \item[\textbf{(5)}] Esto está sacado de la demostración de \parencite{garcia2009hitchin} pero no entiendo bien ni como se define la sección que yo he llamado \(\bm{\zeta}\) y en el paper llamáis \(s_{h,\sigma, \chi}\) ni cuál es la sección holomorfa \(\sigma_{s}\). A partir de ahora sigo vuestra notación en el paper. Tampoco entiendo bien a qué llamáis en el paper \parencite{garcia2009hitchin} \(\barpartial_A\) porque decís que es una estructure holomorfa en el fibrado principal diferenciable subyacente a la reducción \(E_h\), pero citáis el paper de Ignasi \parencite{i1999hitchin}, y en él \(\barpartial_A\) siempre denota un operador de Dolbeault en un fibrado vectorial asociado.
%     \item[\textbf{(6)}] La condición de \(c \in \ii \zenterg \cap \mathfrak{k}\) la he sacado del paper de Ignasi \parencite{i1999hitchin}. En vuestro paper sobre grupo de estructura no conexo, se toma \(c \in \ii \mathfrak{z}_K^\Gamma\) donde \(\mathfrak{z}_K=\mathfrak{z}(\mathfrak{k})\) es el álgebra de Lie del centro de \(K\). En general, no es cierto que \(\zenterg \cap \mathfrak{k}=\mathfrak{z}_K\), aunque sí que tenemos una inclusión \(\zenterg \cap \mathfrak{k}\subset \mathfrak{z}_K\).  Por otro lado, en su artículo \parencite{anchouche2001einstein}, Biswas toma \(c\) directamente en \(\zenterg\), y vosotros en \parencite{garcia2009hitchin}, tomáis \(c \in \mathfrak{z}_\R\) que es el álgebra de Lie del compacto maximal de \(Z\), donde \(Z \subset G\) es el centro de \(G\). Mi pregunta es dónde tengo que tomar el elemento \(c\) en la ecuación. Lo más ``coherente'' con lo que hago para el caso no conexo sería tomar \(c \in \mathfrak{z}_K=\mathfrak{z}(\mathfrak{k})\), pero como en cada paper me parece que se hace una cosa, prefiero preguntar.
%     \item[\textbf{(7)}] He visto que en varios de los artículos habláis de tomar projecciones ortogonales con respecto a la forma bilineal \(\langle \cdot , \cdot \rangle\) en \(\mathfrak{g}\). El problema es que en general, una forma bilineal simétrica compleja en un espacio \(V\) no induce descomposiciones ortogonales ni aún siendo no degenerada. El criterio para que sí que induzca una descomposición ortogonal en un subespacio \(W\) es que su restricción a \(W \times W\) sea no degenerada, pero esto no se sigue de que sea no degenerada en \(V\)  (por ejemplo, en \(\C^2\) la forma bilineal dada en la base canónica por la identidad es no degenerada en todo \(\C^2\), pero no lo es en \(W:=\mathrm{span}_\C(i,1)\), de hecho, la restricción de la forma bilineal en cuestión se anula en \(W\)).
%     \item[\textbf{(8)}] He puesto esta definición de estabilidad buscando ser un poco coherente con las que ponéis en los papers. Los critetios de estabilidad no aparecen escritos exactamente igual en todos los papers, así que no sé si esta estará bien (también he quitado todas las condiciones relativas a la sección \(\varphi\) en el fibrado adjunto, pero esa parte debería estar bien). Más concretamente, he cogido la definición que dais de estabilidad en el paper sin publicar y en el paper \parencite{garcia2009hitchin}. En otros papers como el de Biswas o el Ignasi las condiciones de estabilidad parecen algo diferentes, y no me parece claro de primeras que sean equivalentes. En alguna reunión que hemos tenido me has dicho que sí que lo son, pero aún así preferiría que me confirmaras que esta definición está bien.
%     \item[\textbf{(9)}] El resultado de que \(\mathscr{G}_{k+1}\) tenga una estructura de grupo de Lie de dimensión infinita lo he encontrado solo para \(G\) compacto, conexo, semisimple y matricial. En el paper de Mitter \parencite{mitter1981bundle} dicen que el grupo tiene que ser conexo y compacto aunque yo no he sido capaz de ver donde utilizan la conexión en su prueba (quizás sea alguna propiedad que no sé de la exponencial?). El caso es que lo de que sea un grupo matricial si que me parece muy necesario porque se utiliza muy fuertemente en la construcción. No sé si se suele considerar que todos los grupos con los que se trabaja son matriciales o si hay alguna forma de hacerlo para grupo no matricial, pero yo no la encontré. La verdad es que traté de entender la demostración cuando hice el trabajo y no fui capaz. No tengo nada claro donde necesita todas esas condiciones sobre el grupo de Lie, y no he encontrado más pruebas de que \(\mathscr{G}_k\) sea un grupo de Lie de Hilbert en ningún sitio, así que no sé ni siquiera si es cierto cuando eliminamos todas esas hipótesis y solo pedimos que el grupo sea reductivo.
%     \item[\textbf{(10)}] Lo de que cada órbita contenga una conexión \(\mathscr{C}^\infty\) lo he visto en el articulo de Atiyah-Bott para superficies de Riemann, y para la acción de \(\mathscr{G}^\C\) sobre \(\mathscr{A}^{(1,1)}\). Para la acción de \(\mathscr{G}\) sobre \(\hat\mu^{-1}(0)\) no sé si es cierto, aunque si queremos que el cociente sea el moduli de conexiones de Einstein-Hermite salvo transformación de gauge, tiene sentido. Por otro lado, el otro día hablamos de que todo lo del moment map y el teorema de Kempf-Ness es \textit{lo que motiva} la correspondencia de Hitchin-Kobayashi, o al menos lo que nos lleva a preguntarnos si dicha correspondencia existencia. La parte de la aplicación momento desde luego que veo como se relaciona con el caso finito-dimensional, al fin y al cabo tenemos una aplicación momento y una acción simpléctica. Lo que no veo claro es la parte algebraica. En el kempf-Ness clásico tenemos también nociones de semiestabilidad y estabilidad, pero más allá de que nosotros damos definiciones de estabilidad poliestabilidad etc para fibrados principales, y decidimos darles esos nombres, no veo como esto se tendría que relacionar con la noción finito-dimensional de estabilidad. Es decir, no entiendo por qué uno vería las definiciones que hemos dado de estabilidad, semiestabilidad y poliestabilidad y, si se llamasen de otra forma, pensaría que son análogas a las de dimensión finita (definición \ref{def_stability_fin_dim_KempfNess}).
%     \item[\textbf{(11)}] Que exista un pairing invariante bajo la acción adjunta es algo que afirmáis en el paper pero no dais muchos detalles. He buscado mucho en internet y lo más explícito que he encontrado es una mención a que dichas formas bilineales no degeneradas existes existen en la introducción de este artículo \parencite{Benito2024}. De todas formas, no sé si está bien porque la acción adjunta de \(\hat{G}\) en el centro del álgebra de Lie no tendría que ser trivial en princio si \(\hat{G}\) no es conexo. En resumen, no entiendo de donde sale que exista ese pairing.
%     \item[\textbf{(12)}] No tengo nada claro que eso esté bien, pues de momento no sé si es cierta la relación \(\chi_\zeta(\diff_A \alpha)=\diff \chi_\zeta (\alpha)\). Otro problema que noto es que no sé donde hace falta en nada de lo que hago que al caracter sea antidominante. Sé que tiene que ser necesario en algún momento, no sé si para definir el grado o en algún otro sitio.
%     \item[\textbf{(13)}] Pongamos que tenemos la situación usual \(1 \to G \to \hat{G} \to \Gamma \to 1\). Dado un elemento \(s \in \ii \mathfrak{k}\), se define el grupo parabólico asociado como
%     \[
%     \hat{P}_s:=\corch{g \in \hat{G}: \exp(ts)\ g \ \exp(-ts) \text{ está acotado cuando } t \to \infty }\subset \hat{G}
%     \]
%     y no entiendo con respecto a qué está acotado. Si trabajamos siempre con grupos de matrices, entonces lo vemos en algún \(\mathrm{GL}(n,\C)\), que es un espacio normado y ahí ya tiene perfecto sentido hablar de acotación (y de límites), pero para grupos de Lie abstractos no acabo de ver como lo hacemos. En el texto, lo he interpretado en términos de representaciones unitarias. Esto no lo he visto dicho explícitamente en ningún sitio, pero en \parencite{trautwein2015surveygitpictureyangmills} definen las cosas de un modo un poco diferente, pues piden que el límite exista, no que haya acotación. Para interpretar el límite sí que toman representaciones. No se dice que lo que se hace sea independiente de la representación elegida, cosa que, en caso de ser esta la interpretación correcta, entiendo que debe ser cierta.
%     \item[\textbf{(14)}] Imagino que los subgrupos parabólicos tienen que ser grupos de Lie complejos y subgrupos cerrados de \(G\), pues en los papers tomáis secciones holomorfas de fibrados del tipo \(G/P_s\), de modo que \(G/P_s\) tiene que ser una variedad compleja, y por tanto \(P_s\) debería ser un grupo complejo, pero no he encontrado ningún sitio en que eso esté dicho explícitamente. De hecho, me está constanto mucho encontrar referencias que hablen de subgrupos parabólicos en el contexto de los grupos de Lie. Creo que es cierto que todo grupo reductivo complejo tiene estructure de grupo algebraico, pero en los grupos que tratan con grupos algebraicos como Humprheys \parencite{humphreys1975linear} hacen todo con la topología de Zariski y no me estoy enterando de como pasar estas propiedades a la topología usual. Por ejemplo, para Humprheys \(P \subset G\) es parabólico si y solo si \(G/P\) es una variedad proyectiva \textit{completa}. Completa y compacta en principio no es lo mismo: completa quiere decir, según su propia definición, que para toda variedad proyectiva \(Y\) la segunda proyección \(G/P \times Y \to Y\) es una aplicación cerrada. Según el propio Humprheys, si los espacios son suficientemente buenos y de Hausdorff, y dotamos a todo de la topología producto, entonces esa condición equivale a compacidad, pero él en principio considera todo con la topología de Zariski. Lo mismo en el paper de Ramanathan \parencite{Ramanathan1996} que me recomendaste en su momento, dice que el hecho de que \(P\) sea parabólico es equivalente a que \(G/P\) sea completo, no compacto.
%     \item[\textbf{(15)}] De estas airmaciones hay bastantes cosas de las que no estpy seguro. Casi todo está sacado del remark 3.21 del paper en preparación y de la sección 3.3 del paper de Geometría Dedicata \parencite{barajas2023non}. En primer lugar, que \(G/H\cong \hat{G}/\hat{H}\) no lo tengo muy claro. Si \(\hat{H}\) interseca a todas las componentes conexas de \(\hat{G}\) entonces sí (lo tengo probado en el lema \ref{lemma_HK_non_connected_1}), pero no sé si es cierto. Por otro lado, entiendo que hacemos esto porque queremos que \(G/H\) tenga una acción de \(\hat{G}\), o equivalentemente, una \((\theta,c)\)-twisted \((G, \Gamma)\) action, para poder aplicar así la proposición \ref{prop_gamma_equivariant_bundles} y decir que el fibrado \(P(G/H)\) es \(\Gamma\)-equivariante, y que así tenga sentido hablar de secciones \(\Gamma\)-equivariantes (o \(\Gamma\)-invariantes, que creo que son lo mismo por como definís la \(\Gamma\)-acción, pero creo que no se dice explícitamente).
%     \item[\textbf{(16)}] Lo que pongo en ese remark no lo tengo muy claro. Cuando consideramos \(\hat{G}\) un grupo reductivo, complexificación de \(\hat{K}\), \(G\) es la componente conexa de la identidad de \(\hat{G}\) y definimos \(K:=\hat{K}\cap G\) entiendo que sea compacto, pero no estoy seguro ni de que sea conexo ni de que sea cierto \(K^\C=G\). Tampoco sé si es muy relevante. 
%     \item[\textbf{(17)}] He leído la parte del artículo en que habláis de monodromía y demás, pero no entiendo como deducís que se puede pasar de la ecuación \(\Lambda_\omega F_A=-\ii a\) donde \(a\in \Gamma(P_{\hat{K}} \times_{\hat{K}}\mathfrak{k})\) es una sección con \(\diff_A a=0\) a la misma ecuación con \(\Lambda_\omega F_A=-\ii \alpha\) y \(\alpha \in \mathfrak{z}^\Gamma_K\). 
% \end{enumerate}
% }
